\begingroup
% Local status helpers for the stated thermodynamic Fisher results.
% The importing chapter and appendix retain separate scope groups.
\newcommand{\thermoestablishedstatus}{%
  \unskip\space\status{ESTABLISHED}\par
}
\newcommand{\thermodeclareddefinition}{%
  \unskip\space
  These are declared definitions and conventions; no proof is asserted
  by their adoption. \status{DEFINITION}\par
}
\newcommand{\thermoadoptedhypothesis}{%
  \unskip\space
  These restrictions are adopted hypotheses for the development that
  follows, with the exclusions and scope stated above. \status{HYPOTHESIS}\par
}

\chapter{Proofs for Thermodynamic Fisher Geometry}
\label{app:thermodynamic-fisher-draft-proofs}

This appendix proves the conditional thermodynamic Fisher statements for
the actual finite-free laws and their exact fixed-word pushes. The source,
boundary, support, and regular-chart hypotheses are part of each result;
no unrestricted simultaneous depth and volume limit is asserted.

The earlier spatial appendix supplies the previously stated constrained
specification and local-response setting. The arguments here make the
additional finite-free boundary transfer explicit, then state the exact
inputs used by the retained-label and action calculations. Every regular
inverse is restricted to its declared chart, and every volume limit
retains its boundary and fixed-word convention.\thermoestablishedstatus

\section{Native response and the finite-free boundary bridge}
\label{sec:rg-thermo-native-proof}

The response construction in
\cref{sec:rg-spatial-proof-higher-response,sec:rg-spatial-proof-moving-fisher}
uses finite all-touching specifications and actual infinite-law marginals.
Here the finite outer boundary is free: plaquettes leaving the original
square are omitted. We prove the additional response comparison needed for
these finite laws and their exact fixed-word pushes. Their common limiting
Fisher density does not identify their finite probability laws. In
particular, a free law of a canonical coarse interaction is a fourth
family and does not enter the conclusion below.\thermoestablishedstatus


\subsection{Source box, complete carriers, and fixed words}
\label{sec:rg-thermo-native-families}

\hypothesisheading{Compact native source box}
{hyp:rg-thermo-native-source-box}
Let \(g=(h,j)\), \(b=\tanh h\), and \(t=\tanh j\). Fix a compact
rectangle \(K\) inside
\(h>0\), \(0<j<\operatorname{atanh}(1/12)\), and choose an enlarged
rectangle
\(\widehat K=[h_-,h_+]\times[0,j_+]\), with \(h_->0\) and
\(0<j_+<\operatorname{atanh}(1/12)\), whose positive-\(j\) interior
contains \(K\). Choose \(\tau>0\) and define
\begin{equation}
 q=12\tanh(j_+)e^\tau<1,\qquad
 C=\frac{16q}{1-q},\qquad
 \rho_* =\tanh(h_++2j_+),\qquad d=1-\rho_*>0.
 \label{eq:rg-thermo-native-compact-constants}
\end{equation}
The decay constant \(\tau\) is the constant denoted by \(\theta\) in
\cref{sec:rg-spatial-proof-higher-response}; it is unrelated to the
parameter coordinate \(\theta\) used for a source pullback below. All
constants in this proof are fixed on \(\widehat K\). Its edge \(j=0\)
is used only for a separately labeled one-sided product control.\thermoadoptedhypothesis


\definitionheading{Native free law and its exact pushes}
{def:rg-thermo-native-families}
For \(\Lambda_L=\{0,\ldots,L-1\}^2\), let \(V=L^2\) and retain all
\(2V\) signs \(z_{x,a}\), \(a=1,2\). Put
\begin{equation}
 P_x=z_{x,1}z_{x+e_1,2}z_{x+e_2,1}z_{x,2},\qquad
 r_L(g,z)=\prod_{x\in\Lambda_L}\prod_{a=1}^2(1+bz_{x,a})
          \prod_{x\in\{0,\ldots,L-2\}^2}(1+tP_x).
 \label{eq:rg-thermo-native-free-factor}
\end{equation}
Let \(\nu_{L,g}^{\mathrm f}\) be this raw factor normalized against
counting probability \(4^{-V}\), and let \(Z_L\) be its reference
partition. The unique infinite native Gibbs law is \(\mu_g\), constructed
by its normalized local specifications. For a fixed integer \(B\ge1\),
retain the paths
\(P_B(x,a)=\{(Bx+u e_a,a):0\le u<B\}\) and complete them to a
partition of fine edges with unused singleton blocks. The parameter-independent
map \(C_B\) records their parities. Define
\begin{equation}
 P^{\mathrm f}_{B,L,g}=(C_B)_\#\nu^{\mathrm f}_{BL,g},\qquad
 \mu_{B,g}=(C_B)_\#\mu_g,\qquad
 \mu_{B,L,g}=\mu_{B,g}|_{\Lambda_L}.
 \label{eq:rg-thermo-native-three-families}
\end{equation}
The finite free law, its exact push, and the actual infinite-law marginal
use their own normalized references. A coarse translation by \(x\) is a
fine translation by \(Bx\). Local limits use centered squares, with fine
translation \(-B\lfloor L/2\rfloor\) and coarse translation
\(-\lfloor L/2\rfloor\), or equivalent bulk neighborhoods.

A fixed finite word of integer blocking factors has total linear factor
\(B\). Successive parity products concatenate disjoint straight paths
and equal \(C_B\) on the same finite fine square. Fine area is
\(M=B^2V\) and final coarse area is \(V\). The word, hence \(B\), is
fixed throughout every volume limit. Restricted parity conditionals use
the product of uniform probability references on their prescribed fibers;
unrevealed paths and unused singleton blocks remain unconstrained.\thermodeclareddefinition


\subsection{Constrained contraction and finite-free local response}
\label{sec:rg-thermo-native-free-response}

\lemmaheading{Weighted comparison with omitted outer plaquettes}
{lem:rg-thermo-native-free-contraction}
The constrained weighted comparison and local source-cumulant estimates
of \cref{sec:rg-spatial-proof-constrained,sec:rg-spatial-proof-higher-response}
hold with the constants in \cref{eq:rg-thermo-native-compact-constants}
when plaquettes at the original free outer boundary are omitted. They are
uniform in the finite whole-block inventory, prescribed-or-absent path
pattern, and fixed exterior completion.\thermoestablishedstatus


\paragraph{Proof.}
A plaquette meets four distinct blocks and at most one edge of any retained
path. Conditional on the other fine blocks, a path row is therefore a
positive product sign law, possibly restricted to one parity fiber. To
recall why the row transport constant is independent of path length,
consider one product law conditioned on opposite total parities. It admits
a coupling at Hamming distance exactly one. Induct on its number of bits:
couple the last bits, use identical prefixes when their required parities
agree, and use the induction coupling when those parities differ. In each
branch the conditional prefix marginal is the required product law with
its specified parity, so both complete marginals are correct.

Changing one incident plaquette tilts one coordinate by opposite fields.
If \(u\) collects the positive baseline masses of its two signs, the
coordinate-marginal total variation is
\begin{equation}
 \frac{\sinh(2j)}{\cosh(2u)+\cosh(2j)}\le\tanh j\le\tanh j_+.
 \label{eq:rg-thermo-native-fiber-tilt}
\end{equation}
Couple that bit maximally. A disagreement in a fixed parity fiber needs
one additional discrepancy in the remaining bits, by the preceding
coupling; the mean whole-path Hamming cost is at most \(2\tanh j_+\).
The unconstrained row costs at most \(\tanh j_+\), and a prescribed
one-bit fiber costs zero. Intermediate individual tilts remain positive
product laws on the same support. Composing their coupling kernels and
using the triangle inequality for Hamming distance permits telescoping
all changed incident plaquettes.

For a target block \(i\) and exterior fine edge \(e\), let \(n(i,e)\)
count incident plaquettes. Each fine edge belongs to two plaquettes and
each of those meets three other blocks, so
\(\sum_i n(i,e)\le6\). Omitting outer plaquettes only removes
nonnegative incidences. Finite positive heat baths are irreducible on the
product support; their pair kernel has a stationary law by finite-simplex
Cesaro averaging. If \(D_i\) is its expected block Hamming discrepancy,
canceling the block-update probability gives the inequality in
\cref{eq:rg-spatial-proof-block-hamming} with this restricted incidence
list. There is no extra factor \(B\): \(D_i\) already sums all fine-bit
discrepancies in the updated block.

For target anchor set \(A\), use
\(w_i=e^{-\tau\operatorname{dist}_\infty(\operatorname{anchor}(i),A)}\).
Incident anchors differ by at most one, and consequently
\begin{equation}
 S\le qS+q\sum_{e\in\partial_f D}w_{\operatorname{block}(e)}\delta_e,
 \qquad S=\sum_{i\text{ free}}w_iD_i,
 \quad \delta_e=\mathbf1_{\{\eta_e\ne\zeta_e\}}.
 \label{eq:rg-thermo-native-weighted-comparison}
\end{equation}
This is \cref{eq:rg-spatial-proof-weighted-column} for the actual retained
plaquettes. A shell at anchor distance \(R+1\) from \(A\) has at most
\(8|A|(R+1)\) cells and \(2B^2\) fine edges per cell. Solving
\cref{eq:rg-thermo-native-weighted-comparison} and bounding the probability
of a target discrepancy gives the joint-target bound
\(\min\{1,C|A|B^2(R+1)e^{-\tau(R+1)}\}\).

To obtain group decoupling, condition on all existing free blocks outside
the radius-\(R\) neighborhood of \(A\). Any boundary inherited from the
original inventory, including its omitted outer interactions, is identical
in the two interior conditionals. It produces no near discrepancy. Only
the newly conditioned shell contributes. Each conditional target law is
within the preceding bound of their mixture. Integrating with the exterior
and marginalizing therefore gives, for a vector \(U\) supported on \(A\)
and a vector \(W\) outside its radius-\(R\) neighborhood,
\begin{equation}
 \operatorname{TV}(\mathcal L(U,W),\mathcal L(U)\otimes\mathcal L(W))
 \le\min\{1,C|A|B^2(R+1)e^{-\tau(R+1)}\}.
 \label{eq:rg-thermo-native-group-decoupling}
\end{equation}
The argument permits irregular whole-block inventories. It compares joint
groups, not only pairwise covariances. Infinite constrained laws follow
from consistent finite-cylinder limits as in
\cref{sec:rg-spatial-proof-constrained}; the same bounds pass through
finite-alphabet total variation. No infinite prescribed event is used as
a probability denominator. \hfill\(\square\)

\lemmaheading{Ordered source tails and the finite-free local modulus}
{lem:rg-thermo-native-local-response}
Let a bounded parameter-independent observable \(f\) have anchor support
\(A\), with \(a_0=|A|\ge1\). Its finite-free expectations have actual
compact-uniform derivatives through degree three. If the finite inventory
contains \(A_R=\{x:\operatorname{dist}_\infty(x,A)\le R\}\) and the
finite and infinite constrained patterns agree there, then
\begin{equation}
 \max_{|\alpha|=r}\sup_{g\in\widehat K}
 \left|\partial^\alpha\E_{\nu^D_g}f-
       \partial^\alpha\E_{\mu_g}f\right|
 \le\Delta_r(B,a_0,R)\|f\|_\infty,\qquad 0\le r\le3,
 \label{eq:rg-thermo-native-local-response-limit}
\end{equation}
where \(\mu_g\) denotes the matching constrained law in this statement
and the explicit moduli below tend to zero at fixed \(B,a_0\).\thermoestablishedstatus


\paragraph{Proof.}
Use \(\operatorname{cum}\) for the cumulant of bounded observables. The
moment-partition formula and its constants from
\cref{eq:rg-spatial-proof-response-cumulant} are
\begin{equation}
 \begin{gathered}
 \operatorname{cum}(F_0,\ldots,F_r)
 =\sum_{\pi}(-1)^{|\pi|-1}(|\pi|-1)!
                  \prod_{D\in\pi}\E\prod_{i\in D}F_i,\\
 (k_2,q_2)=(2,6),\qquad (k_3,q_3)=(6,26),\qquad
 (k_4,q_4)=(26,150).
 \end{gathered}
 \label{eq:rg-thermo-native-cumulant-partitions}
\end{equation}
Here \(k_m=\sum_\pi(|\pi|-1)!\) bounds the individual cumulant, and
\(q_m=2\sum_\pi|\pi|!\) multiplies the total-variation change of the
joint argument law. The latter estimate telescopes block moments, each
changing by at most twice total variation. Independent nonempty groups
have zero mixed cumulant because their moment-generating function
factorizes. Expanding its finite logarithm gives the displayed partition
formula, so neither cancellation nor its constants require an infinite
analyticity theorem.

For \(r\in\{1,2,3\}\), let \(T_{i_1,s_1},\ldots,T_{i_r,s_r}\)
be individual edge or plaquette sources for the prescribed ordered
natural derivative directions \(i_1,\ldots,i_r\). Sources may repeat.
Let \(M\) be their maximum base-anchor distance from \(A\). Connect
their labeled anchors and a root representing \(A\) by a minimum spanning
tree. A path to a source at distance \(M\) has at most \(r\) edges and
length at least \(M\), hence a longest tree edge has length at least
\(M/r\). Removing it separates groups with no shorter cross edge.
Plaquette supports enlarge anchor sets by at most one on either side.
For \(M\ge4r\), the groups are therefore separated by a neighborhood
of radius \(\lfloor M/r\rfloor-3\). The root group has at most
\(a_0+4r\) anchors. Apply
\cref{eq:rg-thermo-native-group-decoupling} and cumulant cancellation;
for \(M<4r\), use the individual cumulant bound. With
\begin{equation}
 \begin{split}
 L_r(a_0)&=e^{4\tau}[k_{r+1}+q_{r+1}C(a_0+4r)],\\
 S_r(a_0)&=\frac{(8a_0)^rL_r(a_0)(2r+1)!}
                    {(1-e^{-\tau/r})^{2r+2}},\qquad
 U_r(B,a_0)=S_r(a_0)B^{2r+2},\quad U_0(B,a_0)=1,
 \end{split}
 \label{eq:rg-thermo-native-response-constants}
\end{equation}
each cumulant is bounded by
\(L_r(a_0)B^2(M+1)e^{-\tau M/r}\|f\|_\infty\).
There are at most \(2B^2\) choices of each source type per anchor and at
most \(a_0(2M+1)^2\) anchors within distance \(M\). Thus the number of
ordered tuples of radius at most \(M\) is no greater than
\((8a_0)^rB^{2r}(M+1)^{2r}\). Using this cumulative count also for the
exact-radius count only enlarges the sum. The result is
\begin{equation}
 \begin{split}
 \sum_{s_1,\ldots,s_r}
 |\operatorname{cum}(f,T_{i_1,s_1},\ldots,T_{i_r,s_r})|
 &\le U_r(B,a_0)\|f\|_\infty,\\
 \sum_{\max\operatorname{dist}(x_s,A)>N}
 |\operatorname{cum}(f,T_{i_1,s_1},\ldots,T_{i_r,s_r})|
 &\le U_r(B,a_0)(N+2)^{2r+1}e^{-\tau(N+1)/r}\|f\|_\infty.
 \end{split}
 \label{eq:rg-thermo-native-source-tails}
\end{equation}
For completeness, \((M+1)^k\le k!{M+k\choose k}\) sums the geometric
power series. For its tail write \(M=N+1+n\) and use
\((M+1)^k\le(N+2)^k(n+1)^k\). These are the counts in
\cref{eq:rg-spatial-proof-response-tails}, now valid for the actual
free-boundary source list. Constant exterior-completed sources have zero
higher cumulant; missing outer sources only remove terms.

We now compare local derivatives, without differentiating exterior mixing
weights. For values, condition both laws on a whole-block neighborhood
one interaction layer inside \(A_R\). Every plaquette touching that
smaller region is present in both conditionals. Its exterior assignments
may have different parameter-dependent distributions, but at each fixed
\(g\) the comparison is uniform over both assignments. Integrating it
proves the conservative modulus
\begin{equation}
 \Delta_0(B,a_0,R)=\min\{2,2Ca_0B^2(R+1)e^{-\tau(R-1)}\},
 \qquad R\ge4.
 \label{eq:rg-thermo-native-value-modulus}
\end{equation}
For derivatives, first differentiate each finite normalized law with its
actual source list. Normalized finite exponential differentiation gives
the complete ordered cumulant sum. Put
\(N_R=\lfloor R/2\rfloor-2\). A near source tuple and \(f\) occupy at
most \(a_0+4r\) anchors inside \(A_{N_R+1}\). Their neighborhood of
radius \(R-N_R-2\) has a full extra interaction layer inside \(A_R\).
The joint near-argument total variation is at most
\(C(a_0+4r)B^2(R-N_R+1)e^{-\tau(R-N_R-2)}\).
There are at most \([2B^2a_0(2N_R+1)^2]^r\) near tuples.
Cumulant telescoping bounds their differences; the two far sums are
separately bounded by \cref{eq:rg-thermo-native-source-tails}. Therefore
\begin{equation}
 \begin{split}
 \widetilde\Delta_r(B,a_0,R)={}&
 q_{r+1}[2B^2a_0(2N_R+1)^2]^r
 C(a_0+4r)B^2(R-N_R+1)e^{-\tau(R-N_R-2)}\\
 &+2U_r(B,a_0)(N_R+2)^{2r+1}e^{-\tau(N_R+1)/r},\\
 \Delta_r(B,a_0,R)={}&
 \min\{2U_r(B,a_0),\widetilde\Delta_r(B,a_0,R)\},
 \qquad 1\le r\le3.
 \end{split}
 \label{eq:rg-thermo-native-derivative-moduli}
\end{equation}
For \(R<4\), set \(\Delta_r=2U_r\), including \(r=0\). Each
modulus tends to zero at fixed \(B,a_0\). An omitted original outer
source cannot occur in the common near inventory, so it is covered by the
far comparison rather than erased from a common interior formula.

To justify the infinite derivatives used in this comparison, first exhaust
the matching all-touching fixed-exterior specifications. Each fixed near
tuple has a uniformly convergent finite-alphabet moment law, while
\cref{eq:rg-thermo-native-source-tails} controls both far sums uniformly.
The derivatives through degree three are uniformly Cauchy on compact
rectangles. Their limits are continuous. Pass the finite coordinate
fundamental-theorem-of-calculus identities to the limit, first for the
expectation, then for its first and second derivatives. This identifies
ordinary continuous derivatives through order three, with agreeing mixed
derivatives. The same argument on coordinate segments starting at \(j=0\)
proves continuous one-sided derivatives. The finite-free comparison just
proved converges to those actual derivatives. Neither an infinite
partition nor an infinite-event probability is differentiated.
\hfill\(\square\)

\lemmaheading{Local moving observables}
{lem:rg-thermo-native-moving-response}
For a local \(C^2\) observable \(f_g\), finite and infinite expectations
satisfy
\begin{equation}
 \partial^\alpha\E_g f_g=
 \sum_{\beta\le\alpha}{\alpha\choose\beta}
 \sum_{s_1,\ldots,s_{|\alpha-\beta|}}
 \operatorname{cum}_g(\partial^\beta f_g,T_{i_1,s_1},\ldots,
                                  T_{i_{|\alpha-\beta|},s_{|\alpha-\beta|}}),
 \qquad |\alpha|\le2,
 \label{eq:rg-thermo-native-moving-response}
\end{equation}
where the source directions are those of \(\alpha-\beta\) and an empty
source list means expectation. If all pointwise \(C^2\) jets of \(f\)
on \(A\) are bounded by \(E_f\), then the finite/infinite comparison
of \cref{lem:rg-thermo-native-local-response} gives
\begin{equation}
 \|\E_{\nu^D}f-\E_\mu f\|_{C^2(K)}
 \le E_f[\Delta_0(B,a_0,R)+2\Delta_1(B,a_0,R)+\Delta_2(B,a_0,R)].
 \label{eq:rg-thermo-native-moving-comparison}
\end{equation}\thermoestablishedstatus


\paragraph{Proof.}
Differentiate the finite normalized sum. At each step, a derivative either
acts on \(f\) or inserts the corresponding centered source; successive
insertions give the cumulant. The choices of derivatives on \(f\) give
the multi-index binomial coefficient. Thus the second derivative includes
\(\E f_{ik}\), both source covariance sums involving \(f_i,f_k\),
and the complete three-argument cumulant sum involving \(f\).
This is the finite identity
\cref{eq:rg-spatial-proof-moving-response} on the actual source list.
The source-tail bound permits taking the infinite limit termwise after a
finite cutoff; coordinate FTC identifies the actual derivatives there.
Apply \cref{eq:rg-thermo-native-local-response-limit} separately to the
observable derivatives in the formula. At degree two their total
coefficients are \(1,2,1\), proving the displayed comparison. Both the
observable and its averaging law vary. \hfill\(\square\)

\subsection{Fourth pressure derivatives and the native information tensor}
\label{sec:rg-thermo-native-pressure}

\theoremheading{Native free pressure in \(C^4\)}
{thm:rg-thermo-native-cfour-pressure}
Define the natural statistics and pressure by
\begin{equation}
 \begin{split}
 T_{h,L}&=\sum_{x\in\Lambda_L}\sum_{a=1}^2z_{x,a},\qquad
 T_{j,L}=\sum_{x\in\{0,\ldots,L-2\}^2}P_x,\\
 \mathsf A_L(h,j)&=V^{-1}\log\E_{\mathrm{ref}}
                          e^{hT_{h,L}+jT_{j,L}},\qquad
 \mathfrak p_L=V^{-1}\log Z_L.
 \end{split}
 \label{eq:rg-thermo-native-natural-pressure}
\end{equation}
There are actual \(C^4\) functions \(\mathsf A,\mathfrak p\) on the
positive high-temperature domain such that
\(\mathsf A_L\to\mathsf A\) and
\(\mathfrak p_L\to\mathfrak p\) in \(C^4(K)\) for every compact
source box in \cref{hyp:rg-thermo-native-source-box}. Their exact finite
and limiting conversions are
\begin{equation}
 \begin{split}
 \mathfrak p_L&=\mathsf A_L-2\log\cosh h-q_L\log\cosh j,
       \qquad q_L=\frac{(L-1)^2}{L^2},\\
 \mathfrak p&=\mathsf A-2\log\cosh h-\log\cosh j.
 \end{split}
 \label{eq:rg-thermo-native-pressure-conversion}
\end{equation}\thermoestablishedstatus


\paragraph{Proof.}
For a sign \(z\),
\(1+\tanh(h)z=e^{hz}/\cosh h\). Apply this identity to all \(2V\)
edges and to the \((L-1)^2\) included plaquettes. Both partitions use
normalized counting, so no further cardinality factor enters. This gives
the first line of \cref{eq:rg-thermo-native-pressure-conversion}.

The value limit in \cref{thm:rg-pressure-boundary} applies to this raw
native interaction. Its singleton diameter is zero and each cell has
three plaquette incidences, giving first diameter moment at most
\(\overline M=3t_+/(1-t_+)\). To recall the uniformity, tile a large
square by \(n\)-squares and a leftover set. Terms touching the leftover
set cost its size times a common finite incident norm. Crossing a tile
boundary costs at most \(4\overline M n\) per tile. The exponential
sandwich of finite partition sums therefore traps limiting upper and
lower pressures within \(4\overline M/n\) of \(\mathfrak p_n\).
Their separation is at most \(8\overline M/n\) for every \(n\), so
the limit exists and
\(\|\mathfrak p_L-\mathfrak p\|_{C^0(\widehat K)}
\le4\overline M/L\). The finite conversion gives a uniform value
limit for \(\mathsf A_L\). This step supplies no pressure derivative.

Put \(f_h=z_{0,1}+z_{0,2}\) and \(f_j=P_0\), whose anchor supports
have sizes one and three, and let
\(m_i(g)=\E_{\mu_g}f_i\). These are actual \(C^3\) functions by
\cref{lem:rg-thermo-native-local-response}. Finite differentiation gives
\begin{equation}
 \partial_h\mathsf A_L=\frac1V\sum_{x\in\Lambda_L}
                   \E_{\nu^{\mathrm f}_{L,g}}f_h\circ\vartheta_x,
 \qquad
 \partial_j\mathsf A_L=\frac1V\sum_{x\in\{0,\ldots,L-2\}^2}
                   \E_{\nu^{\mathrm f}_{L,g}}f_j\circ\vartheta_x,
 \label{eq:rg-thermo-native-pressure-gradient}
\end{equation}
where \(\vartheta_x\) translates the local observable to \(x\).
For \(0\le r\le3\), define
\begin{equation}
 M_r=\max\{2U_r(1,1),U_r(1,3)\},\qquad
 D_r(R)=\max\{2\Delta_r(1,1,R),\Delta_r(1,3,R)\}.
 \label{eq:rg-thermo-native-gradient-bounds}
\end{equation}
For \(L\ge2(R+2)\), at most \(4(R+2)L\) base cells lack a complete
target support and its radius-\(R\) neighborhood. The bulk differences
of local derivatives are bounded by \(D_r(R)\). Each finite and infinite
boundary derivative is bounded by \(M_r\); an omitted plaquette is an
identically zero finite term. Therefore
\begin{equation}
 \max_{i\in\{h,j\},|\alpha|=r}
 \|\partial^\alpha\partial_i\mathsf A_L-
                       \partial^\alpha m_i\|_{C^0(K)}
 \le D_r(R)+8M_r\frac{R+2}{L},\qquad 0\le r\le3.
 \label{eq:rg-thermo-native-pressure-gradient-boundary}
\end{equation}
At fixed \(R\), take \(L\) large; then take \(R\) large. Equivalently,
for a prescribed tolerance first choose \(R\) for its local modulus and
then \(L\) for its boundary fraction. The gradients converge in \(C^3\).

On any coordinate segment in the enlarged rectangle,
\begin{equation}
 \mathsf A_L(g+u e_i)-\mathsf A_L(g)
       =\int_0^u\partial_i\mathsf A_L(g+v e_i)dv.
 \label{eq:rg-thermo-native-pressure-ftc}
\end{equation}
Uniform value and gradient convergence pass through the integral, giving
\(\partial_i\mathsf A=m_i\). Apply the same argument to the first,
second and third gradient derivatives in
\cref{eq:rg-thermo-native-pressure-gradient-boundary}. It proves actual
\(C^4\) regularity, convergence of every derivative through order four,
and equality of mixed derivatives. Since \(|q_L-1|\le2/L\) and
\(q_L\) is parameter independent, the conversion proves the same result
for raw pressure. In particular, its fourth derivative requires only
three local-response derivatives: its first derivative is already a local
source density. The argument never differentiates a surface log-ratio
estimate. \hfill\(\square\)

\propositionheading{Native Fisher density and a source pullback}
{prop:rg-thermo-native-ctwo-metric}
The actual native free family satisfies
\begin{equation}
 \frac1V I(\nu^{\mathrm f}_{L,g})=\nabla_g^2\mathsf A_L(g)
 \longrightarrow H(g):=\nabla_g^2\mathsf A(g)
 \quad\text{in }C^2(K),
 \qquad
 H=\nabla_g^2\mathfrak p+
      \operatorname{diag}(2\operatorname{sech}^2h,
                              \operatorname{sech}^2j).
 \label{eq:rg-thermo-native-fisher-pressure}
\end{equation}
If \(\chi:\Theta\subset\R^d\to\R^2\),
\(\chi(\theta)=(h(\theta),j(\theta))\), is \(C^4\) and sends a
neighborhood of a compact parameter set into an admitted source box, then
\begin{equation}
 g_L(\theta)=\frac1V I(\nu^{\mathrm f}_{L,\chi(\theta)})
 =D\chi(\theta)^\top\nabla_g^2\mathsf A_L(\chi(\theta))D\chi(\theta)
 \longrightarrow
 g(\theta)=D\chi(\theta)^\top H(\chi(\theta))D\chi(\theta)
 \quad\text{in }C^2.
 \label{eq:rg-thermo-native-source-pullback}
\end{equation}\thermoestablishedstatus


\paragraph{Proof.}
The natural normalized log score is
\((T_{h,L},T_{j,L})-\E(T_{h,L},T_{j,L})\). Its covariance is exactly
\(V\nabla_g^2\mathsf A_L\). The pressure theorem and its scalar
conversion yield \cref{eq:rg-thermo-native-fisher-pressure}. The ordinary
chain rule for the normalized log probability multiplies the source score
by \(D\chi^\top\), proving the finite pullback. Its \(C^2\) convergence
follows from the \(C^2\) source-Hessian limit and the product and
composition rules. In contrast,
\begin{equation}
 \nabla_\theta^2(\mathsf A_L\circ\chi)
 =D\chi^\top(\nabla_g^2\mathsf A_L)D\chi
   +\sum_{i=1}^2(\partial_i\mathsf A_L)\nabla_\theta^2\chi_i.
 \label{eq:rg-thermo-native-nonnatural-hessian}
\end{equation}
The last term is absent from the score covariance. Thus neither the raw
pressure Hessian nor a Hessian after nonlinear source substitution can be
used without its stated correction. \hfill\(\square\)

\propositionheading{A native finite-energy information floor}
{prop:rg-thermo-native-finite-energy-floor}
On an admitted compact source box put
\(H_*=h_++2j_+\), \(m_*=(1+e^{2H_*})^{-1}\), and
\(v_* =\operatorname{sech}^2 H_*\). For \(L\ge12\),
\begin{equation}
 \frac1V I(\nu^{\mathrm f}_{L,g})\succeq
       \frac{v_*m_*^6}{8}I_2,
 \qquad H(g)\succeq\frac{v_*m_*^6}{8}I_2.
 \label{eq:rg-thermo-native-compact-floor}
\end{equation}
Consequently the pullback in
\cref{eq:rg-thermo-native-source-pullback} has rank
\(\operatorname{rank}D\chi\); positivity on an arbitrary \(d\)-parameter
chart is not presumed.\thermoestablishedstatus


\paragraph{Proof.}
Every edge conditioned on all other signs has field magnitude at most
\(H_*\), so either sign has conditional probability at least \(m_*\)
and conditional variance at least \(v_*\). Choose horizontal edges based
at \((1+3u,1+3v)\), where \(u,v\ge0\) and both coordinates are at most
\(L-2\). There are \(\lfloor L/3\rfloor^2\ge V/16\) such edges for
\(L\ge12\), and no two share a plaquette. Condition on all unselected
signs. The selected signs are then independent. For
\(F=w_hT_{h,L}+w_jT_{j,L}\), a selected edge \(e\) has coefficient
\(w_h+w_jR_e\), where \(R_e\) is the sum of the two neighboring
three-sign products. Thus
\begin{equation}
 \operatorname{Var}(F)\ge
 \E\operatorname{Var}(F\given\text{unselected signs})
 \ge v_*\sum_{e\text{ selected}}\E(w_h+w_jR_e)^2.
 \label{eq:rg-thermo-native-separated-variance}
\end{equation}
Each selected edge has six distinct neighboring signs, all unselected and
inside the square. Setting all six positive forces \(R_e=2\). Setting
the horizontal edge immediately above and the horizontal edge immediately
below negative, and the four vertical neighbors positive, forces
\(R_e=-2\). Each of these disjoint assignments has probability at least
\(m_*^6\). To see this under partial conditioning, reveal the six signs
successively: averaging the full one-edge conditional lower bound over
all unspecified signs preserves \(m_*\) at every step. Therefore
\begin{equation}
 \frac{\operatorname{Var}(F)}V
 \ge\frac{v_*m_*^6}{16}
            [(w_h+2w_j)^2+(w_h-2w_j)^2]
 =\frac{v_*m_*^6}{8}(w_h^2+4w_j^2)
 \ge\frac{v_*m_*^6}{8}|w|^2.
 \label{eq:rg-thermo-native-directional-floor}
\end{equation}
This is the finite matrix bound, and the proved source-Hessian limit
inherits it. A positive definite middle matrix in \(D\chi^\top H D\chi\)
has kernel exactly \(\ker D\chi\), proving the rank statement. On a
compact full-column-rank chart the smallest singular value of \(D\chi\)
has a positive minimum. No positive metric on more than two independent
source coordinates, or global quotient of a singular pullback, follows.
\hfill\(\square\)

\subsection{Predictive rows under the finite free law}
\label{sec:rg-thermo-native-predictive-rows}

\lemmaheading{Finite-free predictive rows through third order}
{lem:rg-thermo-native-row-jets}
Let \(D\) be a finite anchor inventory, \(x\in D\), and \(S\subset D\)
a set of cells preceding \(x\) in the fixed lexicographic order. Write
\(p^D_S(g,s;y)\) for the four-state row of the exact finite free push,
conditional on \(Y_S=y_S\). It is positive and \(C^3\), uniformly in
\(S,y,D,x\) with the bounds below. The corresponding infinite-law row
\(p_S\) prescribes exactly this past and leaves every future path free.\thermoestablishedstatus


\paragraph{Proof.}
The finite conditioning restricts precisely the observed path fibers and
retains every original internal plaquette. It deletes no additional
interaction, and an absent future is an unconstrained future. A currently
free target path conditioned on all other fine blocks has independent
signs with fields \(h+j a_e\), \(|a_e|\le2\). Its parity mean is
\(G_g=\prod_{e\in P}\tanh(h+j a_e)\), of magnitude at most
\(\rho_*^B\), with support in at most nine anchors. Successively revealing
the two target parities gives
\begin{equation}
 p^D_S(g,s;y)\ge\frac{(1-\rho_*^B)^2}{4}\ge\frac{d^2}{4}.
 \label{eq:rg-thermo-native-row-positivity}
\end{equation}
Finite constrained response applies at the original free boundary by
\cref{lem:rg-thermo-native-free-contraction}.

For \(n\le3\), assigning \(n\) labeled derivatives to the \(B\)
hyperbolic-tangent factors gives
\(\|\partial^\beta G_g\|_\infty\le g_nB^n\rho_*^{B-n}\),
where \(g_n=2^n n!\) and \(|\beta|=n\). The first three tangent
derivatives are bounded by \(1,2,6\), each \(j\) derivative contributes
at most two, and the product of the factorials of assigned derivative
orders is at most \(n!\). At \(B<n\) the negative exponent only enlarges
the bound. Apply the moving-observable formula to \(\E G_g\), keeping
both derivatives of the amplitude and source insertions. Exactly as in
\cref{eq:rg-spatial-proof-mean-jets}, this yields
\begin{equation}
 \begin{split}
 D_n&=g_n+\sum_{a=0}^{n-1}{n\choose a}g_aS_{n-a}(9),\\
 M_0(B)&=\rho_*^B,\qquad
 M_n(B)=D_nB^{2n+2}\rho_*^{B-n},\qquad 1\le n\le3.
 \end{split}
 \label{eq:rg-thermo-native-parity-jets}
\end{equation}
These bound every order-\(n\) target parity-mean derivative. The actual
infinite derivatives and their one-sided continuation follow from
\cref{lem:rg-thermo-native-local-response}.

For one Bernoulli factor \((1+s m)/2\), the denominator \(1+s m\) is
at least \(d\). Direct logarithmic differentiation gives
\begin{equation}
 A_1=\frac{M_1}{d},\qquad
 A_2=\frac{M_2}{d}+\frac{M_1^2}{d^2},\qquad
 A_3=\frac{M_3}{d}+\frac{3M_2M_1}{d^2}+\frac{2M_1^3}{d^3}.
 \label{eq:rg-thermo-native-log-amplitudes}
\end{equation}
A whole-cell log row \(l=\log p^D_S\) has order-\(n\) derivatives
bounded by \(\ell_n=2A_n\). Its second target-path factor conditions
on the same already revealed first parity, retaining their dependence.
These are the finite-free counterparts of
\cref{eq:rg-spatial-proof-log-jet-amplitudes}.

Define the row-Fisher observable
\(\mathsf I^D_{S,ik}=\sum_s p^D_S l_i l_k\), whose configuration
argument depends only on \(y_S\). Its derivatives at fixed \(y\) are
\begin{equation}
 \begin{split}
 \partial_u\mathsf I^D_{S,ik}
 &=\sum_s p^D_S(l_ul_il_k+l_{iu}l_k+l_il_{ku}),\\
 \partial_v\partial_u\mathsf I^D_{S,ik}
 &=\sum_s p^D_S\bigl[(l_{uv}+l_ul_v)l_il_k
 +l_u(l_{iv}l_k+l_il_{kv})+l_v(l_{iu}l_k+l_il_{ku})\\
 &+l_{iuv}l_k+l_{iu}l_{kv}+l_{iv}l_{ku}+l_il_{kuv}\bigr].
 \end{split}
 \label{eq:rg-thermo-native-row-fisher-jets}
\end{equation}
In particular \(\partial_u p=p l_u\) and
\(\partial_v\partial_u p=p(l_{uv}+l_ul_v)\) have been included.
The ten expanded second-derivative monomials give the common entrywise
bound
\begin{equation}
 E_B^F=\max\{\ell_1^2,\ell_1^3+2\ell_1\ell_2,
          \ell_1^4+5\ell_1^2\ell_2+2\ell_1\ell_3+2\ell_2^2\},
 \qquad H_B^F=\max(1,\ell_1,\ell_2,\ell_3),
 \label{eq:rg-thermo-native-row-fisher-bound}
\end{equation}
for the observable and its first two derivatives.

Suppose two target parity means have jet differences bounded by
\(v_0,v_1,v_2,v_3\). Subtract the differentiated logarithms and use
\(|a^{-k}-b^{-k}|\le k|a-b|/d^{k+1}\). The one-factor log-jet differences
are bounded by
\begin{equation}
 \begin{split}
 \mathcal T_1(v)&=v_1/d+M_1v_0/d^2,\\
 \mathcal T_2(v)&=v_2/d+(M_2v_0+2M_1v_1)/d^2+2M_1^2v_0/d^3,\\
 \mathcal T_3(v)&=v_3/d+(M_3v_0+3M_2v_1+3M_1v_2)/d^2\\
 &\quad +(6M_2M_1v_0+6M_1^2v_1)/d^3+6M_1^3v_0/d^4.
 \end{split}
 \label{eq:rg-thermo-native-log-moduli}
\end{equation}
These reconstruct \cref{eq:rg-spatial-proof-log-jet-moduli}. A complete
row uses \(2\mathcal T_n\); its probability difference is at most
\(v_0\), because each of its two Bernoulli factors changes by at most
\(v_0/2\). Telescope the at most five factors in each monomial of
\cref{eq:rg-thermo-native-row-fisher-jets}, over four states, to obtain
\begin{equation}
 \mathcal W_B(v)=\min\{2E_B^F,
 200(H_B^F)^4[v_0+2\mathcal T_1(v)+2\mathcal T_2(v)+2\mathcal T_3(v)]\}.
 \label{eq:rg-thermo-native-row-modulus-map}
\end{equation}
It bounds every entrywise row-Fisher derivative difference through degree
two, with a separately valid cap \(2E_B^F\).

For two past patterns in the same finite \(D\) agreeing through \(Q_R\),
let \(a_R=\lfloor R/2\rfloor-1\) for \(R\ge2\), and put
\(N_n=2^nq_{n+1}C(1+4n)\). Their parity-jet differences obey
\begin{equation}
 \begin{split}
 V_0(B,R)&=\min\{2\rho_*^B,2CB^2(R+1)e^{-\tau(R+1)}\},\\
 V_n(B,R)&=\min\bigl\{2M_n(B),B^{2n+2}\bigl[
 N_n(2a_R+1)^{2n}(R-a_R)e^{-\tau(R-a_R)}\\
 &+2S_n(1)(a_R+2)^{2n+1}e^{-\tau(a_R+1)/n}\bigr]\bigr\},
 \qquad 1\le n\le3.
 \end{split}
 \label{eq:rg-thermo-native-past-jet-modulus}
\end{equation}
For \(R=0,1\), take \(V_n=2M_n\). To justify the same finite-domain
modulus near an original free boundary, split source tuples at \(a_R\).
Near joint supports occupy at most \(1+4n\) anchors in \(Q_{a_R+1}\);
there are at most \([2B^2(2a_R+1)^2]^n\) tuples. Group comparison on
the remaining buffer gives the first term. Original omitted interactions
are identical in the two laws and produce no near discrepancy. Each far
sum obeys \cref{eq:rg-thermo-native-source-tails}. This is the argument
of \cref{eq:rg-spatial-proof-jet-modulus} on the retained finite source
list, including irregular intersections with the original inventory.
Consequently
\begin{equation}
 \|\mathsf I^D_S-\mathsf I^D_{S'}\|_{C^2(K),\infty,\mathrm{entry}}
 \le\omega_B(R):=\mathcal W_B(V_0,V_1,V_2,V_3).
 \label{eq:rg-thermo-native-past-fisher-modulus}
\end{equation}
At fixed \(B\), \(\omega_B(R)\) is bounded by a finite constant times
\((R+1)^7e^{-\tau R/6}\). The order-three response supplies the exponent
\(\tau/6\); a first-order rate has not been substituted.

If the target has \(Q_R\subset D\), compare a finite row with the
infinite row of the same local past. Apply
\cref{eq:rg-thermo-native-local-response-limit} with \(a_0=1\) separately
to each target parity. When comparing the second factor, the first parity
is prescribed identically on both sides. Thus
\begin{equation}
 \|\mathsf I^D_S-\mathsf I_{S'}\|_{C^2(K),\infty,\mathrm{entry}}
 \le\Omega_B(R):=\mathcal W_B\bigl(
 \Delta_0(B,1,R),\Delta_1(B,1,R),\Delta_2(B,1,R),\Delta_3(B,1,R)\bigr)
 \longrightarrow0.
 \label{eq:rg-thermo-native-free-row-comparison}
\end{equation}
This bound is pointwise in the common past coordinates and controls actual
row derivatives. It has not yet differentiated their expectations under
a moving law. \hfill\(\square\)

\lemmaheading{Weighted shells under the moving unconditioned law}
{lem:rg-thermo-native-moving-shells}
For a local \(C^2\) observable supported in \(Q_N\), put
\(a_N=(2N+1)^2\) and
\begin{equation}
 \|f\|_{2,K,\infty}=\max_{|\alpha|\le2}\sup_{g\in K,y}
                                  |\partial^\alpha f(g,y)|,
 \qquad W_B(N)=1+2S_1(a_N)B^4+S_2(a_N)B^6.
 \label{eq:rg-thermo-native-shell-response-weight}
\end{equation}
The original finite free and infinite unconstrained fine laws satisfy
\begin{equation}
 \|\E_gf_g\|_{C^2(K)}\le W_B(N)\|f\|_{2,K,\infty}.
 \label{eq:rg-thermo-native-shell-expectation-bound}
\end{equation}
Every finite-past row expectation, including targets at the finite free
boundary, has entrywise \(C^2\) bound
\begin{equation}
 Z_B^F=W_B(0)E_B^F+\sum_{n\ge1}W_B(n)\omega_B(n-1)<\infty,
 \qquad
 \mathcal T_B(N)=\sum_{n>N}W_B(n)\omega_B(n-1)\longrightarrow0.
 \label{eq:rg-thermo-native-shell-tail}
\end{equation}
Here \(\mathcal T_B(N)\) bounds the \(C^2\) error of truncating the
past at \(Q_N\). The same statements hold for the infinite full-past
row and its actual expectation under \(\mu_g\).\thermoestablishedstatus


\paragraph{Proof.}
Pull the coarse observable back along the fixed parity map. Its fine
support uses at most \(a_N\) anchors. In the first derivative of its
expectation, the moving-response formula has an expectation of \(f_i\)
and a source covariance sum. In the second derivative it has an
expectation of \(f_{ik}\), the two source covariance sums of \(f_i,f_k\),
and the two-source cumulant sum of \(f\). Applying
\cref{eq:rg-thermo-native-source-tails} to these terms proves
\cref{eq:rg-thermo-native-shell-expectation-bound}. The infinite version
follows by finite-near and uniform-tail passage before coordinate FTC,
as in \cref{lem:rg-thermo-native-moving-response}. This supplies the
finite-free counterpart of
\cref{eq:rg-spatial-proof-varying-expectation-bound}.

For a finite past \(S\), let \(S_n=S\cap Q_n\),
\(f_n^D=\mathsf I^D_{S_n}\),
\(d_0^D=f_0^D\), and \(d_n^D=f_n^D-f_{n-1}^D\). The two prefixes
agree through \(Q_{n-1}\), so
\begin{equation}
 \|d_0^D\|_{2,K,\infty}\le E_B^F,\qquad
 \|d_n^D\|_{2,K,\infty}\le\omega_B(n-1),\qquad n\ge1.
 \label{eq:rg-thermo-native-shell-differences}
\end{equation}
Each difference is a function only of \(y_{S_n}\), although its
parameter-dependent coefficients were computed from the entire finite
law. Its pointwise derivatives include that coefficient dependence.
Thus configuration locality is exactly what the moving-response bound
needs; no derivative of the finite normalizer is discarded.

The degree of \(S_1(a)\) in \(a\) is at most two, and that of
\(S_2(a)\) at most three. Hence \(W_B(n)\) grows at most as a degree-six
polynomial in \(n\), whereas \(\omega_B(n-1)\) is bounded by a degree-seven
polynomial times \(e^{-\tau n/6}\) at fixed \(B\). Their product is
summable. The finite series terminates exactly once \(Q_n\) contains
\(S\). The common weighted sum bounds every finite expected row,
independently of its target, domain and remote past shape, and bounds its
truncation remainder by \(\mathcal T_B(N)\).

For an infinite past the pointwise row jets are uniformly Cauchy by
\cref{eq:rg-thermo-native-past-fisher-modulus}. Their shell expectations
and derivatives through degree two are uniformly Cauchy under the same
majorant. Passing finite-sum coordinate FTC identities identifies an
actual \(C^2\) expectation and allows termwise differentiation. Absolute
convergence permits reordering the source and shell sums. Throughout,
the averaging law is the original unconditioned fine law. The constrained
law defines the row; it is not substituted for this outer expectation.
\hfill\(\square\)

\subsection{Exact finite-free pushes and the three-cutoff limit}
\label{sec:rg-thermo-native-free-push}

\theoremheading{Exact pushed-free Fisher density}
{thm:rg-thermo-native-free-push-ctwo}
For every fixed integer \(B\ge1\) and compact positive source box,
\begin{equation}
 \frac1{L^2}I(P^{\mathrm f}_{B,L,g})
 \longrightarrow\mathcal I_B(g):=
          \E_{\mu_g}\mathsf I_{\mathcal P}(g,Y)
 \quad\text{in }C^2(K),
 \label{eq:rg-thermo-native-free-push-limit}
\end{equation}
where \(\mathcal P\) is the complete lexicographic past and every future
path is unconstrained. This is the same density as the actual
infinite-law marginal limit in
\cref{eq:rg-spatial-proof-actual-fisher-limit}. The finite laws in the two
limits need not agree.\thermoestablishedstatus


\paragraph{Proof.}
Order the cells of \(\Lambda_L\) lexicographically and use the true
finite predictive rows of \(P^{\mathrm f}_{B,L,g}\). Their log scores add
to the complete score. Each row score has conditional mean zero given
its finite past, since differentiating its finite normalization gives
zero. An earlier row score is measurable with respect to any later past;
conditioning on that past therefore kills every cross expectation. With
\(S_{x,L}=\Lambda_L\cap\{y:y<x\}\), the exact finite identity is
\begin{equation}
 I(P^{\mathrm f}_{B,L,g})=
    \sum_{x\in\Lambda_L}\E_{\nu^{\mathrm f}_{BL,g}}
                              \mathsf I^{D_L}_{S_{x,L}}(g,Y).
 \label{eq:rg-thermo-native-finite-score-rows}
\end{equation}
It is an equality of finite \(C^2\) functions, so its first two derivatives
are still exact finite identities.

Fix an inner past cutoff \(N\ge0\), then an outer radius
\(R\ge N+4\). Consider a bulk target \(x\) with
\(x+Q_R\subset\Lambda_L\) and translate it to zero. Its finite past
and the complete past have exactly the same intersection with \(Q_N\).
Truncating each of their expected row observables to that common prefix
costs at most \(2\mathcal T_B(N)\) in entrywise \(C^2\), by
\cref{lem:rg-thermo-native-moving-shells}.

The finite and infinite truncated rows are functions of the same finite
past coordinates. Their pointwise \(C^2\) difference is at most
\(\Omega_B(R)\), by
\cref{eq:rg-thermo-native-free-row-comparison}. Its expectation under the
moving finite law is therefore bounded by \(W_B(N)\Omega_B(R)\).
This step differentiates the finite averaging law; it does not multiply a
pointwise error by an uncontrolled extensive score.

It remains to compare expectations of the same infinite truncated row
under \(\nu^{\mathrm f}_{BL,g}\) and \(\mu_g\). Its pointwise jets are
bounded by \(E_B^F\), its support is \(Q_N\), and this support has a
radius-\((R-N)\) neighborhood in \(D_L\). Thus
\cref{eq:rg-thermo-native-moving-comparison} bounds the remaining cost by
\(E_B^F[\Delta_0+2\Delta_1+\Delta_2](B,a_N,R-N)\).
Combining these three comparisons gives the bulk bound
\begin{equation}
 \begin{split}
 \mathcal B_B(N,R)={}&2\mathcal T_B(N)+W_B(N)\Omega_B(R)\\
 &+E_B^F[\Delta_0(B,a_N,R-N)+2\Delta_1(B,a_N,R-N)
                                  +\Delta_2(B,a_N,R-N)].
 \end{split}
 \label{eq:rg-thermo-native-three-cutoff-bulk}
\end{equation}
At most \(4RL\) cells are outside the bulk for \(L\ge2R\). Each finite
and infinite expected-row jet is bounded by \(Z_B^F\), including near the
free outer boundary. Stationarity identifies the infinite target
expectation at each bulk translation. Dividing
\cref{eq:rg-thermo-native-finite-score-rows} by \(L^2\) therefore yields
\begin{equation}
 \left\|L^{-2}I(P^{\mathrm f}_{B,L,g})-\mathcal I_B(g)
                                      \right\|_{C^2(K),\mathrm{entry}}
 \le\mathcal B_B(N,R)+8Z_B^F\frac RL.
 \label{eq:rg-thermo-native-three-cutoff-bound}
\end{equation}
Given a tolerance, first choose \(N\) for the weighted shell tail. Next
choose \(R\) for the two local comparisons at that fixed \(N\), and
finally choose \(L\) for the boundary fraction. This order retains the
growth of the support-dependent constants in \(a_N\); no bound uniform
in \(N\) has been presumed. All terms vanish in that order, with \(B\)
fixed. This proves the asserted \(C^2\) limit.

For actual infinite-law marginals the same finite score-row identity uses
\(\mu_g\) as its outer expectation at every volume. Their finite pasts
and full past agree through \(Q_R\) at all but \(4RL\) sites. The weighted
shell argument then gives
\(2\mathcal T_B(R)+8Z_B^F R/L\), exactly the proof of
\cref{eq:rg-spatial-proof-actual-fisher-limit}. Hence its limit is the
same \(\mathcal I_B\). At \(B=1\), \(C_1\) is the identity, and
\cref{eq:rg-thermo-native-fisher-pressure,eq:rg-thermo-native-free-push-limit}
are limits of the same finite matrices. Consequently
\begin{equation}
 \mathcal I_1=H=\nabla_g^2\mathsf A.
 \label{eq:rg-thermo-native-identity-density}
\end{equation}
For \(B>1\), an exact push need not have a two-statistic exponential
representation, so it has not been identified with a Hessian of native
or canonical coarse pressure. \hfill\(\square\)

\propositionheading{Area normalization of the fixed-word comparison}
{prop:rg-thermo-native-word-area}
For fixed positive integers \(a,d_0\),
\begin{equation}
 \mathcal I_{ad_0}\preceq d_0^2\mathcal I_a,
 \qquad
 \frac{\mathcal I_{ad_0}}{(ad_0)^2}\preceq\frac{\mathcal I_a}{a^2},
 \qquad \mathcal I_B\preceq B^2H.
 \label{eq:rg-thermo-native-word-area}
\end{equation}\thermoestablishedstatus


\paragraph{Proof.}
The laws \(P^{\mathrm f}_{a,d_0L}\) and
\(P^{\mathrm f}_{ad_0,L}\) are obtained from the same fine free square
of side \(ad_0L\), and the latter is the exact \(d_0\)-path push of the
former. Differentiating a finite parameter-independent channel gives an
output score equal to the conditional expectation of the input score.
Conditional covariance decomposition proves
\(I(P^{\mathrm f}_{ad_0,L})\preceq I(P^{\mathrm f}_{a,d_0L})\).
Divide by the final coarse area \(L^2\) and take the two fixed-depth
limits to obtain the first inequality; the second is its rearrangement.
Direct comparison with the fine native law gives the last inequality.
These are tensor inequalities at each source point. Differentiating a
positive-semidefinite inequality does not order its derivatives.
\hfill\(\square\)

\subsection{Product controls and nonempty regular charts}
\label{sec:rg-thermo-native-regular-charts}

\propositionheading{Product-boundary scores and fixed-depth positive charts}
{prop:rg-thermo-native-fixed-depth-chart}
For every prescribed finite list of depths, the tensors \(\mathcal I_B\)
have a common nonempty open positive-\(j\) source region on which each is
positive definite. Each compact subchart has a positive minimum eigenvalue
at every listed depth. The minimum need not persist over all integer
depths.\thermoestablishedstatus


\paragraph{Proof.}
First consider the native product control \(j=0\). Write
\(b=\tanh h\) and \(v=1-b^2\). Independent fine signs give
\begin{equation}
 H(h,0)=\begin{pmatrix}
 2v&4b^3v\\
 4b^3v&1-b^8+4b^6v
 \end{pmatrix},\qquad
 \det H(h,0)=2v^3(1+2b^2+3b^4)>0.
 \label{eq:rg-thermo-native-product-hessian}
\end{equation}
The field covariance counts two signs per cell. The mixed covariance
counts four edges per plaquette, each contributing \(b^3v\). A plaquette
has variance \(1-b^8\) and four neighbors sharing one edge, each with
covariance \(b^6v\); all other plaquette covariances vanish. Expanding
the determinant gives
\(2v(1-b^8)-8b^6v^2=2v^3(1+2b^2+3b^4)\).
These are derivatives of the full source family at the boundary, including
its transverse plaquette response. They are not obtained by holding
\(j\) identically zero while differentiating. The one-sided response
construction justifies the values and their continuation. This product
argument is separate from the positive-coupling compact finite-energy
bound in \cref{prop:rg-thermo-native-finite-energy-floor}.

For a general fixed \(B\), prescribe any complete lexicographic past,
leave future paths free, and work at \(j=0\). The constrained fine law
is a product of its block-fiber laws. A free target path has parity mean
\(b^B\). For its parity value \(s\), the contrast of any one path edge is
\begin{equation}
 a(s)=\E[z_e\given Y=s]-b
     =\frac{s b^{B-1}(1-b^2)}{1+s b^B},\qquad
 \E a=0,\qquad
 v_B=\E a^2=\frac{b^{2B-2}(1-b^2)^2}{1-b^{2B}}>0.
 \label{eq:rg-thermo-native-parity-contrast}
\end{equation}
Multiply the independent sign law by its parity indicator to obtain this
conditional mean, and sum its two parity outcomes to obtain the variance.
Conditional variance decomposition also gives \(v_B\le1-b^2\le1\).
The two target contrasts \(a_1,a_2\) are independent and centered.
Finite source differentiation gives the actual field row score
\(s_h=B(a_1+a_2)\).

The plaquette at the common starting vertex is the unique plaquette
meeting both outgoing target paths. Its other two edges are future edges
at \(B=1\) and unused singleton edges at \(B>1\), hence both have mean
\(b\). Its centered contribution is
\(b^2a_1a_2+b^3(a_1+a_2)\). Each other incident plaquette meets one
target path, and contributes a multiple of that path's contrast with
coefficient of magnitude at most one. There are \(2B-1\) such incidences
per target path. Nonincident plaquettes cancel by block independence.
Thus the actual plaquette score is
\begin{equation}
 s_j=b^2a_1a_2+\alpha_1a_1+\alpha_2a_2,
 \qquad |\alpha_i|\le2B.
 \label{eq:rg-thermo-native-product-row-scores}
\end{equation}
These finite source identities pass to the one-sided limiting derivatives
by \cref{lem:rg-thermo-native-local-response}, agreeing with the explicit
construction in \cref{sec:rg-spatial-proof-regular-chart}.

The bilinear contrast is orthogonal to both linear contrasts. Therefore
\begin{equation}
 \begin{split}
 I_{hh}^{\mathrm{row}}&=2B^2v_B,\qquad
 I_{hj}^{\mathrm{row}}=Bv_B(\alpha_1+\alpha_2),\\
 I_{jj}^{\mathrm{row}}&=v_B(\alpha_1^2+\alpha_2^2)+b^4v_B^2,\\
 \det I^{\mathrm{row}}&=B^2v_B^2(\alpha_1-\alpha_2)^2
                                      +2B^2b^4v_B^3,\qquad
 \operatorname{Tr}I^{\mathrm{row}}\le11B^2v_B.
 \end{split}
 \label{eq:rg-thermo-native-product-row-floor}
\end{equation}
The smaller eigenvalue is at least the determinant divided by the trace,
hence at least \(2b^4v_B^2/11\), uniformly in the past. On a compact
positive field interval, \(b\in[b_-,b_+]\subset(0,1)\), define
\begin{equation}
 \eta_{0,B}=\frac2{11}b_-^{4B}(1-b_+^2)^4>0,\qquad
 Q_B=16(A_1^3+A_1A_2)>0,\qquad
 \delta_B=\min\{j_+/2,\eta_{0,B}/(2Q_B)\}>0.
 \label{eq:rg-thermo-native-chart-constants}
\end{equation}
The floor follows from
\(v_B\ge b_-^{2B-2}(1-b_+^2)^2\). The row-score vector has norm at
most \(2\sqrt2 A_1\), and its \(j\) derivative has norm at most
\(2\sqrt2 A_2\). Differentiating its four-state Fisher sum gives matrix
norm at most \(16A_1^3\) for the probability derivative and at most
\(16A_1A_2\) for the two score derivatives. Its total derivative is
therefore at most \(Q_B\). The actual one-sided derivatives and FTC give
\begin{equation}
 I^{\mathrm{row}}(h,j\given\text{past})\succeq\frac{\eta_{0,B}}2I_2,
 \qquad \mathcal I_B(h,j)\succeq\frac{\eta_{0,B}}2I_2,
 \qquad 0<j<\delta_B.
 \label{eq:rg-thermo-native-positive-chart}
\end{equation}
The second inequality is the expectation of the first. Enlarge the field
interval before fixing the constants to obtain an open neighborhood of
the desired compact interval. For a prescribed finite depth list, take
finite minima of widths and floors. There is no positive all-depth minimum
in this argument. \hfill\(\square\)

\propositionheading{Geometry on the admitted regular charts}
{prop:rg-thermo-native-regular-geometry}
On the compact regular charts just exhibited, the normalized finite-free
and exact-push Fisher metrics converge with their inverses in \(C^2\),
their Levi-Civita coefficients in \(C^1\), and their Riemann, Ricci and
scalar curvatures in \(C^0\). The actual-marginal metrics obey the same
conclusion. A general source pullback needs an admitted full-column-rank
chart and its positive compact singular-value minimum.\thermoestablishedstatus


\paragraph{Proof.}
Write \(G\) for a limiting positive metric and \(G_L\) for its indicated
normalized finite metric. The proved \(C^2\) convergence makes
\(G_L\) positive for every sufficiently large \(L\), with a common
inverse bound. Differentiating \(GG^{-1}=I\) gives
\begin{equation}
 \partial_iG^{-1}=-G^{-1}(\partial_iG)G^{-1},
 \label{eq:rg-thermo-native-inverse-derivative}
\end{equation}
and a second differentiation is a finite sum of products of inverse and
metric derivatives through order two. These formulas prove \(C^2\)
inverse convergence. In local coordinates the formulas
\begin{equation}
 \begin{split}
 \Gamma^l_{ik}&=\frac12G^{lm}
  (\partial_iG_{mk}+\partial_kG_{mi}-\partial_mG_{ik}),\\
 R^l{}_{kiv}&=\partial_i\Gamma^l_{vk}-\partial_v\Gamma^l_{ik}
       +\Gamma^l_{im}\Gamma^m_{vk}-\Gamma^l_{vm}\Gamma^m_{ik}
 \end{split}
 \label{eq:rg-thermo-native-connection-curvature}
\end{equation}
then give the asserted derivative orders. Their torsion-free,
metric-compatible and tensorial interpretation is the one proved in
\cref{eq:rg-spatial-proof-fisher-connection-curvature}; only metric
derivatives through degree two occur. Ricci and scalar contractions use
the same convergent inverse. These statements concern the normalized
metrics, whose constant volume rescaling cannot be suppressed when
comparing scalar curvature.

These source coordinates also give local charts of the actual law. At a
point, choose a sufficiently large actual finite marginal with positive
Fisher matrix. Its positive probability vector satisfies
\(I=(Dp)^\top\operatorname{diag}(1/p)Dp\), so two atom probabilities
have a nonsingular two-by-two Jacobian. Multiply that two-coordinate map
by the inverse Jacobian at the point and shrink a convex ball until its
derivative differs from the identity by less than \(1/2\). Segment
integration gives a lower Lipschitz bound. For nearby target values the
iteration \(x\mapsto x+y-F(x)\) preserves a smaller ball and contracts
distances by \(1/2\); its successive differences form a convergent
geometric series, yielding a local inverse. Differentiating its defining
identity gives the inverse derivatives. Thus distinct nearby parameters
give distinct actual laws. This does not assert a global quotient of a
singular tensor. Fisher curvature, gauge-connection curvature and base
curvature retain their different underlying spaces. \hfill\(\square\)

\subsection{Why scalar boundary control is insufficient}
\label{sec:rg-thermo-native-boundary-obstruction}

\propositionheading{A bounded boundary ratio with divergent information density}
{prop:rg-thermo-native-boundary-score-counterexample}
Uniform subvolume control of log-density ratios and agreement of bulk
cylinder limits do not imply convergence of volume-normalized Fisher
tensors, even for positive smooth finite families on one fixed carrier
at each volume.\thermoestablishedstatus


\paragraph{Proof.}
Use the complete square carrier with \(2V\) signs, \(V=L^2\). Let
\(P_L\) be the fair product law, independent of a real parameter
\(\theta\). Choose one sign on the boundary of the recentered square
and let \(Q_{L,\theta}\) change only its mean to
\(m_L(\theta)=a\sin(V\theta)\), with fixed \(0<a<1/2\).
All probabilities are positive and smooth on the same open parameter
interval, and for every configuration
\begin{equation}
 \left|\log\frac{Q_{L,\theta}}{P_L}\right|
 \le\max\{-\log(1-a),\log(1+a)\}.
 \label{eq:rg-thermo-native-boundary-ratio-control}
\end{equation}
The varying sign leaves every fixed recentered cylinder for sufficiently
large \(L\), so the two sequences have the same bulk-cylinder limit.
For a sign of mean \(m\), differentiating its two probabilities gives
Fisher information \((m')^2/(1-m^2)\). Independence of all other signs
therefore gives
\begin{equation}
 I(Q_{L,\theta})=
   \frac{a^2V^2\cos^2(V\theta)}{1-a^2\sin^2(V\theta)},
 \qquad I(Q_{L,0})=a^2V^2,
 \qquad V^{-1}I(Q_{L,0})=a^2V\longrightarrow\infty,
 \label{eq:rg-thermo-native-boundary-score-growth}
\end{equation}
whereas \(I(P_L)=0\). The witness violates the native uniform local jet
bounds: its boundary response grows with \(V\). It refutes the scalar
boundary inference, not the native theorems above. \hfill\(\square\)

At the separately labeled native product control \(j=0\), the raw
partition is one for every \(h\), so
\(\partial_h^2\mathfrak p(h,0)=0\), while
\(H_{hh}(h,0)=2\operatorname{sech}^2h>0\). Thus even in the native
family the raw-pressure Hessian needs the conversion in
\cref{eq:rg-thermo-native-fisher-pressure}. More generally,
\(f_L(\theta)=L^{-1}\sin(L^2\theta)\) converges uniformly to zero but
has derivative \(L\) at zero, showing directly why the pressure proof
must use local response rather than differentiation of its value error.\thermoestablishedstatus


An exact finite shifted-source identity permits the preceding native and
fixed-word proofs to be applied after its admitted \(C^4\) source map,
with scalar evidence accounted for separately in the raw geometry weights.
A general positive likelihood or a free canonical coarse-interaction
family requires its own response comparison or an exact finite family
identity. The common infinite law and a surface log-ratio alone do not
supply that derivative bridge. No arbitrary-observation, critical,
continuum, simultaneous depth-volume, or all-depth inverse conclusion is
included in this proof fragment.\thermoestablishedstatus

% Host: existing Theory macros; no preamble or chapter heading.
% Source transcription: saved P4--P6, corrected LG01--LG07,
% independent label IR01--IR19, and IF02--IF04.
% Native C4 pressure, C2 exact-push, and positive-chart inputs are linked

\section{Proofs for retained labels, ties, and reciprocal-volume experiments}
\label{sec:rg-thermo-label-proof}

This section uses the actual normalized finite experiments of
\Cref{sec:rg-thermo-finite-experiment}. The sample space retains the
geometry label $\gamma\in\mathfrak G$, even when conditional branches
coincide. All supports are fixed on the parameter neighborhood in use;
structural zero atoms are omitted from score sums and have no conditional
score. The limit theorems below state their branch response inputs explicitly.
For the admitted native families these are supplied by
\Cref{sec:rg-thermo-native-proof}; other branch families require their
own matching input. The label conclusions also require the stated
representative growth, gaps, or actual center ratios and do not follow
from pressure convergence alone. Every volume limit is taken at fixed
catalog and fixed block word.\thermoestablishedstatus


\subsection{Native inputs and effective representatives}
\label{sec:rg-thermo-label-inputs}

\definitionheading{Volume families and ordinary derivative norms}
{def:rg-thermo-label-family}
Let $\Theta\subset\R^d$ be open, let $V=L^2$, and let
$p_{\gamma,L,\theta}$ be a positive normalized $C^3$ mass family on its
fixed active finite carrier. Carriers may depend on $L$. Define
\begin{equation}
 P_{L,\theta}(\gamma,x)=W_{\gamma,L}(\theta)p_{\gamma,L,\theta}(x),
 \qquad
 W_{\gamma,L}=\frac{e^{\lambda_{\gamma,L}}}
 {\sum_\eta e^{\lambda_{\eta,L}}},
 \qquad
 \xi_{\gamma,L}=\nabla\lambda_{\gamma,L}.
 \label{eq:rg-thermo-label-family}
\end{equation}
For a compact $K$ inside an open neighborhood, $\|f\|_{C^k(K)}$ denotes
the maximum of the suprema of its ordinary coordinate derivatives
through degree $k$, with a fixed finite-dimensional norm for vector and
matrix entries. Derivatives of order zero are included. Put
$g_{\gamma,L}=V^{-1}I(p_{\gamma,L})$ and
$g_L=V^{-1}I(P_L)$. Unless stated otherwise, every convergence in this
section is uniform in the indicated $C^k$ norm on each such compact set.
The bulk function $s_\gamma$ below is an effective log-mass density; it
is distinct from the physical scalar in a raw-factor representation.
\status{DEFINITION}

\propositionheading{Conditional map from native response to branch inputs}
{prop:rg-thermo-label-native-map}
Consider the complete native free square with normalized counting
reference $4^{-V}$, all $2V$ edge signs, and $(L-1)^2$ internal
plaquettes. Write
\begin{equation}
 T_L=\sum_e z_e,\qquad U_L=\sum_{p\ \mathrm{internal}}P_p,\qquad
 \mathsf A_L(h,j)=\frac1V\log\sum_x4^{-V}
                         e^{hT_L(x)+jU_L(x)}.
 \label{eq:rg-thermo-label-natural-pressure}
\end{equation}
The actual $C^4$ convergence $\mathsf A_L\to\mathsf A$ is supplied by
\Cref{thm:rg-thermo-native-cfour-pressure} on the compact source boxes
of \Cref{hyp:rg-thermo-native-source-box}, inside
$h>0$, $0<j<\operatorname{arctanh}(1/12)$. This uses the finite free
law of \Cref{def:rg-thermo-native-families}, with the same complete
carrier and normalized reference as above. If
$\chi_\gamma:\Theta\subset\R^d\to\R^2$,
$\chi_\gamma(\theta)=(h_\gamma(\theta),j_\gamma(\theta))$, is $C^4$
on a neighborhood of $K$ and sends that neighborhood into a common
compact admitted source box, then
\Cref{prop:rg-thermo-native-ctwo-metric} gives
\begin{align}
 \mathfrak p_{\gamma,L}
 &:=\frac1V\log Z_{\gamma,L}
 =\mathsf A_L\circ\chi_\gamma-2\log\cosh h_\gamma
   -\frac{(L-1)^2}{V}\log\cosh j_\gamma,
 \label{eq:rg-thermo-label-raw-pressure}\\
 g_{\gamma,L}
 &=(D\chi_\gamma)^\top
       (D^2\mathsf A_L\circ\chi_\gamma)D\chi_\gamma
 \longrightarrow
 g_\gamma=(D\chi_\gamma)^\top
       (D^2\mathsf A\circ\chi_\gamma)D\chi_\gamma
 \quad\hbox{in }C^2(K).
 \label{eq:rg-thermo-label-native-pullback}
\end{align}
In particular $\mathfrak p_{\gamma,L}$ has an actual $C^4$ limit
$\mathfrak p_\gamma$ and the branch densities have bounded $C^2$ norms.
This records the application of the native derivative theorem; its
local-response and finite-free boundary proof is
\Cref{thm:rg-thermo-native-cfour-pressure}.\thermoestablishedstatus


\paragraph{Proof.}
For a sign $z$,
$1+(\tanh h)z=e^{hz}/\cosh h$. Apply this identity to every edge and
plaquette, with $b=\tanh h$ and $t=\tanh j$, and count the factors.
This gives \cref{eq:rg-thermo-label-raw-pressure} exactly. The normalized
branch mass is the exponential family with sufficient statistic
$(T_L,U_L)$: its scalar hyperbolic-cosine factors cancel in its own
normalization. Differentiating the finite sum in
\cref{eq:rg-thermo-label-natural-pressure} gives
$D^2\mathsf A_L=V^{-1}\operatorname{Cov}(T_L,U_L)$ in natural source
coordinates. Multiplying the centered natural score by
$(D\chi_\gamma)^\top$ gives the displayed tensor pullback. Ordinary
chain and product rules through degree two use derivatives of
$\mathsf A_L$ through degree four and of $\chi_\gamma$ through degree
three. The declared $C^4$ source maps more than supply these derivatives.
The raw-pressure convergence follows by composition and by
$(L-1)^2/V\to1$. No derivative of a $C^0$ pressure error is taken.
The Hessian of $\mathsf A_L\circ\chi_\gamma$ would additionally contain
$\sum_a(\partial_a\mathsf A_L\circ\chi_\gamma)D^2\chi_{\gamma,a}$;
the Hessian of raw pressure would also retain the scalar correction.
Neither is the Fisher pullback in general. $\square$

For an exact fixed-word push, use the distinct family in
\Cref{def:rg-thermo-native-families}. Let $C_B$ be its fixed
parameter-independent straight-path parity channel, including the
unrevealed paths and unused singleton blocks in the fine carrier.
A fixed finite word of the admitted integer path maps concatenates to
this $C_B$ on the same fine square, as in
\Cref{def:rg-thermo-native-families,prop:rg-thermo-native-word-area}.
Write $p_{\gamma,BL}^{\rm free}C_B$ for the actual push from $B^2V$
fine cells to $V=L^2$ coarse cells. The source-coordinate theorem
\Cref{thm:rg-thermo-native-free-push-ctwo}, followed by the same $C^4$
source map, supplies
\begin{equation}
 \frac1V I_\theta(p_{\gamma,BL}^{\rm free}C_B)
 \longrightarrow (D\chi_\gamma)^\top
      (\mathcal I_B\circ\chi_\gamma)D\chi_\gamma
 \quad\hbox{in }C^2(K).
 \label{eq:rg-thermo-label-push-input}
\end{equation}
Here $\mathcal I_B$ is the natural-source Fisher density per coarse cell
of actual infinite-law marginals. To make the parameter dimension
explicit, at each finite volume the normalized score chain rule gives
\begin{equation}
 I_\theta(p_{\gamma,BL}^{\rm free}C_B)
 =(D\chi_\gamma)^\top
 I_{(h,j)}(P^{\rm f}_{B,L,\chi_\gamma(\theta)})D\chi_\gamma.
 \label{eq:rg-thermo-label-push-pullback}
\end{equation}
The source theorem is a $C^2$ tensor limit in two natural coordinates;
composition with $\chi_\gamma$ and the ordinary product rule through
degree two give \cref{eq:rg-thermo-label-push-input} in the declared
$d$ coordinates. No rank assumption is needed for this semidefinite
pullback. The theorem's three-cutoff boundary estimate
\cref{eq:rg-thermo-native-three-cutoff-bound} identifies its actual
finite-free limit with the actual-marginal density while keeping these
finite families distinct. The actual-marginal theorem alone, a scalar
pressure estimate, an arbitrary channel, or a separately normalized
free canonical coarse family does not supply this input. The word and
$B$ remain fixed before the volume limit.\thermoestablishedstatus

\propositionheading{Effective growth in both normalization modes}
{prop:rg-thermo-label-representatives}
Suppose the positive raw weights have actual $C^3$ logarithms and satisfy
\begin{equation}
 \|\log a_{\gamma,L}\|_{C^3(K)}\leq C_KV,
 \qquad V^{-1}\log a_{\gamma,L}\longrightarrow\alpha_\gamma
                  \quad\hbox{in }C^3(K).
 \label{eq:rg-thermo-label-raw-growth}
\end{equation}
Using \Cref{thm:rg-thermo-native-cfour-pressure} through
\Cref{prop:rg-thermo-label-native-map}, valid effective representatives
and bulk limits for the latent experiment are
\begin{align}
 \lambda_{\gamma,L}^{\rm f}
 &=\log a_{\gamma,L}+\log Z_{\gamma,L},
 &s_\gamma^{\rm f}&=\alpha_\gamma+\mathfrak p_\gamma,\label{eq:rg-thermo-label-factor-growth}\\
 \lambda_{\gamma,L}^{\rm h}
 &=\log a_{\gamma,L},
 &s_\gamma^{\rm h}&=\alpha_\gamma.\label{eq:rg-thermo-label-hierarchy-growth}
\end{align}
Each satisfies
\begin{equation}
 \|\lambda_{\gamma,L}\|_{C^3(K)}\leq C_KV,
 \qquad \lambda_{\gamma,L}/V\longrightarrow s_\gamma
                                      \quad\hbox{in }C^3(K).
 \label{eq:rg-thermo-label-effective-growth}
\end{equation}
The same assertion holds for the admitted all-edge selected posterior,
with the effective representatives given below. Common normalizers must
be removed algebraically before imposing this growth condition.\thermoestablishedstatus


\paragraph{Proof.}
The finite formulas in
\cref{eq:rg-thermo-finite-log-masses,eq:rg-thermo-finite-labeled-fisher}
give \cref{eq:rg-thermo-label-factor-growth,eq:rg-thermo-label-hierarchy-growth}.
Replacing every $\lambda_{\gamma,L}$ by
$\lambda_{\gamma,L}+c_L(\theta)$ leaves every $W_{\gamma,L}$ unchanged
and adds the same $\nabla c_L$ to every $\xi_{\gamma,L}$. Their covariance
is unchanged as a function, including all its available derivatives.
Thus a hierarchy generated from raw $a$ uses $\log a$ for
\cref{eq:rg-thermo-label-effective-growth}. One need not, and cannot in
general, impose the same bound on
$\log W_\gamma=\log a_\gamma-\log\sum_\eta a_\eta$ at a tie.
For example, with $a_0=1$, $a_1=e^{V\theta}$, the second derivative of
$\log(a_0+a_1)$ at zero is $V^2/4$. Its exact common cancellation
preserves the stated raw-weight domain. A fixed positive $C^3$
prescribed hierarchy $w_\gamma$ is a separate admissible case with
$\lambda_{\gamma,L}=\log w_\gamma$ and bulk limit zero.

For the all-edge binary sensor of
\cref{eq:rg-thermo-finite-sensor-shift}, let its mean parameter be
$\eta_\gamma(\theta)$ and set
\begin{equation}
 A_\gamma^{\rm sens}=\frac{1+b_\gamma\eta_\gamma}{2},\qquad
 b'_\gamma=\frac{b_\gamma+\eta_\gamma}{1+b_\gamma\eta_\gamma}.
 \label{eq:rg-thermo-label-sensor-data}
\end{equation}
Assume these data are $C^4$ on a neighborhood of $K$, the sensor is
strictly positive, and both original $(h_\gamma,j_\gamma)$ and shifted
$(\operatorname{arctanh}b'_\gamma,j_\gamma)$ sources lie in common
compact admitted boxes. Multiplication over exactly $2V$ signs gives
$Z_\gamma^o=(A_\gamma^{\rm sens})^{2V}Z_{\gamma,L}(b'_\gamma,t_\gamma)$.
The selected conditional branch is exactly the shifted native free law.
Up to a common log mass, the posterior representatives are
\begin{align}
 \lambda_{\gamma,L}^{o,\rm f}
 &=\log a_{\gamma,L}+2V\log A_\gamma^{\rm sens}
                 +\log Z_{\gamma,L}(b'_\gamma,t_\gamma),
 \label{eq:rg-thermo-label-posterior-factor}\\
 \lambda_{\gamma,L}^{o,\rm h}
 &=\log a_{\gamma,L}+2V\log A_\gamma^{\rm sens}
       +\log Z_{\gamma,L}(b'_\gamma,t_\gamma)
       -\log Z_{\gamma,L}(b_\gamma,t_\gamma).
 \label{eq:rg-thermo-label-posterior-hierarchy}
\end{align}
Divide by $V$ and apply \Cref{thm:rg-thermo-native-cfour-pressure} to
both the original and shifted sources in their admitted boxes. The
bulk limits are respectively
$\alpha_\gamma+2\log A_\gamma^{\rm sens}+\mathfrak p'_\gamma$ and
$\alpha_\gamma+2\log A_\gamma^{\rm sens}
 +\mathfrak p'_\gamma-\mathfrak p_\gamma$ in $C^3$.
\Cref{prop:rg-thermo-native-ctwo-metric} applied to the shifted source
supplies the conditional branch input; its exact fixed-word push uses
\Cref{thm:rg-thermo-native-free-push-ctwo} and
\cref{eq:rg-thermo-label-push-pullback}. All sensor derivatives remain
in these expressions. For a prescribed
hierarchy replace $\log a$ by $\log w$ and keep the same physical
partition difference. Strict positivity makes the finite selected event
well defined. This argument is confined to this exact source-shift
identity; it asserts no derivative theorem for an arbitrary positive
likelihood or a null infinite observation event. $\square$

When the volume family instead denotes a fixed-word push, its branches
in \cref{eq:rg-thermo-label-family} use coarse volume $V$, whereas its
label weights are inherited from the fine volume $M=B^2V$. Thus
$\lambda_{\gamma,BL}/V\to B^2s_\gamma$ in $C^3$ and the same strict
winner is selected. This conversion is required before applying a
volume-normalized statement. It neither changes $B$ with $L$ nor turns
an arbitrary stochastic channel into spatial blocking.

\subsection{Strict phases and regular charts}
\label{sec:rg-thermo-label-strict}

\theoremheading{Strict-phase convergence with both metric derivatives}
{thm:rg-thermo-label-strict}
Assume \cref{eq:rg-thermo-label-effective-growth} on a neighborhood of
$K$ and actual branch convergence $g_{\gamma,L}\to g_\gamma$ in
$C^2(K)$ for every label. Suppose a fixed label $\gamma_*$ obeys
$s_{\gamma_*}-s_\gamma\geq\delta_K>0$ on $K$ for every
$\gamma\ne\gamma_*$. Then
\begin{equation}
 \left\|g_L-g_{\gamma_*}\right\|_{C^2(K)}\longrightarrow0.
 \label{eq:rg-thermo-label-strict-limit}
\end{equation}
More precisely, for $j=0,1,2$ the losing label term obeys
\begin{equation}
 \left\|D^j\left[V^{-1}
       \operatorname{Cov}_{W_L}(\xi_{\gamma,L})\right]\right\|_{K}
 \leq C_K V^{j+1}e^{-\delta_KV/2},
 \label{eq:rg-thermo-label-strict-covariance-bound}
\end{equation}
and the weighted branch-density error relative to its winning branch
has $C^2(K)$ norm at most $C_KV^2e^{-\delta_KV/2}$ for all sufficiently
large $L$. Only three derivatives of the effective log masses are used.\thermoestablishedstatus


\paragraph{Proof.}
Uniform convergence of $\lambda/V$ implies
$\lambda_{\gamma,L}-\lambda_{\gamma_*,L}\leq-\delta_KV/2$ for each
loser when $L$ is large. The softmax denominator includes the winning
mass, so $W_{\gamma,L}\leq e^{-\delta_KV/2}$. Write
$\overline\xi_i=\sum_\eta W_\eta\xi_{\eta,i}$, suppressing $L$.
Finite differentiation gives
\begin{align}
 \partial_iW_\gamma&=W_\gamma(\xi_{\gamma,i}-\overline\xi_i),
 \label{eq:rg-thermo-label-weight-first}\\
 \partial_j\overline\xi_i
 &=\sum_\eta W_\eta\partial_j\xi_{\eta,i}
                  +\operatorname{Cov}_W(\xi_i,\xi_j),
 \label{eq:rg-thermo-label-mean-derivative}\\
 \partial_{ij}W_\gamma
 &=W_\gamma\big[(\xi_{\gamma,i}-\overline\xi_i)
                  (\xi_{\gamma,j}-\overline\xi_j)
             +\partial_j\xi_{\gamma,i}-\partial_j\overline\xi_i\big].
 \label{eq:rg-thermo-label-weight-second}
\end{align}
Every $\xi$ and its first two derivatives is $O(V)$. The covariance in
\cref{eq:rg-thermo-label-mean-derivative} is $O(V^2)$, so for every
ordinary multi-index $\beta$ of degree at most two,
\begin{equation}
 \|\partial^\beta W_{\gamma,L}\|_K
 \leq C_K V^{|\beta|}e^{-\delta_KV/2}
       \quad(\gamma\ne\gamma_*).
 \label{eq:rg-thermo-label-losing-derivatives}
\end{equation}
The same bound applies to $W_{\gamma_*,L}-1$ and its derivatives,
since it is minus the sum of losing weights.

Put $d_\gamma=\xi_\gamma-\xi_{\gamma_*}$ and
$m_d=\sum_{\gamma\ne\gamma_*}W_\gamma d_\gamma$. Exact centering gives
\begin{equation}
 \operatorname{Cov}_W(\xi)
 =\sum_{\gamma\ne\gamma_*}W_\gamma d_\gamma d_\gamma^\top
                                      -m_dm_d^\top.
 \label{eq:rg-thermo-label-strict-centered-covariance}
\end{equation}
Each derivative of $d_\gamma$ through degree two is $O(V)$. A term
with $j$ total derivatives in the first sum has bound
$CV^{j+2}e^{-\delta_KV/2}$: at most $j$ derivatives land on a weight,
and the two score factors each contribute at most $V$.
The product $m_dm_d^\top$ has the same polynomial bound and two
exponential factors, hence also this bound. Dividing by $V$ proves
\cref{eq:rg-thermo-label-strict-covariance-bound}.

For the branch term use the exact identity
\begin{equation}
 \sum_\gamma W_\gamma g_{\gamma,L}-g_{\gamma_*,L}
 =\sum_{\gamma\ne\gamma_*}W_\gamma
                         (g_{\gamma,L}-g_{\gamma_*,L}).
 \label{eq:rg-thermo-label-strict-branch-error}
\end{equation}
The branch differences have uniformly bounded $C^2$ norms. The ordinary
Leibniz sum, including all derivatives of $W$, and
\cref{eq:rg-thermo-label-losing-derivatives} yield its stated $C^2$
bound. Add the winning branch's $C^2$ limit error and the covariance
error, using \cref{eq:rg-thermo-finite-labeled-fisher}. Exponentials
dominate all the displayed powers of $V$, proving
\cref{eq:rg-thermo-label-strict-limit}. $\square$

The native positivity input is
\Cref{prop:rg-thermo-native-finite-energy-floor}. On its compact source
box use $H_*=h_++2j_+$, $m_*=(1+e^{2H_*})^{-1}$, and
$v_*=\operatorname{sech}^2H_*$. Its exact finite and limiting bounds are
\begin{equation}
 D^2\mathsf A_L\succeq\frac{v_*m_*^6}{8}\operatorname{Id}_2
       \quad(L\geq12),\qquad
 D^2\mathsf A\succeq\frac{v_*m_*^6}{8}\operatorname{Id}_2.
 \label{eq:rg-thermo-label-native-floor-input}
\end{equation}
The separated-edge conditional-variance proof is given there. The
pullback in \cref{eq:rg-thermo-label-native-pullback} therefore has
kernel $\ker D\chi_\gamma$ and rank $\operatorname{rank}D\chi_\gamma$,
since a positive matrix has zero quadratic form only at the zero vector.\thermoestablishedstatus
\propositionheading{Inverse, connection, and curvature on regular charts}
{prop:rg-thermo-label-regular-geometry}
Let symmetric tensor fields $G_L\to G$ in $C^2$ on a neighborhood of a
compact parameter chart, and suppose $G$ is positive definite there.
Then $G_L^{-1}\to G^{-1}$ in $C^2$, the Levi-Civita coefficients converge
in $C^1$, and the Riemann, Ricci, and scalar curvatures converge in
$C^0$ on the compact chart. This applies to the strict-phase density
of \cref{thm:rg-thermo-label-strict} precisely on such a chart.\thermoestablishedstatus


\paragraph{Proof.}
For the fixed Euclidean coordinate norm, minimize
$z^\top G(\theta)z$ on the compact product of the chart with the unit
sphere. Positivity gives a minimum $\epsilon>0$. Uniform convergence
implies $G_L\succeq(\epsilon/2)\operatorname{Id}$ for large $L$.
The identity $G_L^{-1}-G^{-1}=G_L^{-1}(G-G_L)G^{-1}$ proves uniform
inverse convergence. Direct differentiation of $G_LG_L^{-1}=\operatorname{Id}$
gives, writing $H_L=G_L^{-1}$,
\begin{align}
 \partial_iH_L&=-H_L(\partial_iG_L)H_L,\\
 \partial_{ij}H_L
 &=H_L(\partial_jG_L)H_L(\partial_iG_L)H_L
  +H_L(\partial_iG_L)H_L(\partial_jG_L)H_L
  -H_L(\partial_{ij}G_L)H_L.
 \label{eq:rg-thermo-label-inverse-derivatives}
\end{align}
Every factor converges uniformly, proving $C^2$ inverse convergence.
Use the explicit coefficients and curvature convention
\begin{align}
 \Gamma^k_{ij}(G_L)
 &=\frac12(G_L)^{k\ell}
   (\partial_i(G_L)_{j\ell}+\partial_j(G_L)_{i\ell}
                                      -\partial_\ell(G_L)_{ij}),\\
 R^\ell{}_{kij}(G_L)
 &=\partial_i\Gamma^\ell_{jk}-\partial_j\Gamma^\ell_{ik}
    +\Gamma^\ell_{im}\Gamma^m_{jk}
    -\Gamma^\ell_{jm}\Gamma^m_{ik}.
 \label{eq:rg-thermo-label-connection-curvature}
\end{align}
The first line and its first derivatives converge; the second line and
its contractions therefore converge continuously. Lowering the first
index multiplies by the convergent metric, and scalar contraction uses
the convergent inverse. These arguments prove exactly the stated
orders, without third derivatives of the metric. $\square$

For a parameter-independent positive constant $V$, replacing $G_L$ by
$I_L=VG_L$ leaves the connection, type-$(1,3)$ Riemann tensor, and
covariant Ricci tensor unchanged. The fully covariant Riemann tensor
is multiplied by $V$, and scalar curvature is divided by $V$, as follows
by substitution in \cref{eq:rg-thermo-label-connection-curvature} and its
contractions. A finite positive Fisher tensor implies that the
Jacobian of the finite probability map has full column rank: its Gram
form is the Fisher quadratic form with positive diagonal weights.
A full-rank probability minor supplies a locally identifiable chart.
It does not give global identifiability.

The native floor makes the natural two-source chart nonempty. A general
$d$-parameter native pullback has rank at most two; directions acting
only on losing labels can disappear in the strict-phase limit. For the exact push input \cref{eq:rg-thermo-label-push-input},
\Cref{prop:rg-thermo-native-fixed-depth-chart} supplies the required
source floor on its specific positive-$j$ region. On its compact
positive field interval, the constants $\eta_{0,B}>0$ and $\delta_B>0$
in \cref{eq:rg-thermo-native-chart-constants} give
$\mathcal I_B(h,j)\succeq(\eta_{0,B}/2)\operatorname{Id}_2$ for
$0<j<\delta_B$, by \cref{eq:rg-thermo-native-positive-chart}.
For a prescribed finite list $\mathcal B$ choose a parameter compact
whose source image lies in this field interval and
$0<j<\min_{B\in\mathcal B}\delta_B$. The source derivative on the
chosen chart must have full column rank, with compact minimum singular
value $\sigma_K>0$. Its pulled-back limiting metrics then have the
common floor
$(\sigma_K^2/2)\min_{B\in\mathcal B}\eta_{0,B}>0$.
Together with the strict-phase convergence above this supplies the
hypotheses of \Cref{prop:rg-thermo-label-regular-geometry}, consistently
with the branch geometry in
\Cref{prop:rg-thermo-native-regular-geometry}. An arbitrary positive
high-temperature source image is not thereby shown to lie in these
small-$j$ charts; that inclusion and full column rank remain explicit
conditions for this inverse statement. No positive floor uniform in
unbounded depth, inverse on a singular stratum, or quotient inferred
solely from constant rank is asserted.

\subsection{Fixed ties and lower information scales}
\label{sec:rg-thermo-label-ties}

\propositionheading{Squared-volume Fisher tensor at an actual finite-weight tie}
{prop:rg-thermo-label-tie-scale}
Fix $\theta_0$ and assume \cref{eq:rg-thermo-label-effective-growth} on
a neighborhood. Let $S$ be the nonempty maximizing set of
$s_\gamma(\theta_0)$, suppose the actual finite weights satisfy
$W_{\gamma,L}(\theta_0)\to w_\gamma$, and assume
$I(p_{\gamma,L})(\theta_0)=O(V)$ for every label. Write
$v_\gamma=\nabla s_\gamma(\theta_0)$ and
$\overline v=\sum_\gamma w_\gamma v_\gamma$. Then $w$ is supported on
$S$ and
\begin{equation}
 \frac{I(P_L)(\theta_0)}{V^2}
 \longrightarrow C_w(v)
 :=\sum_\gamma w_\gamma(v_\gamma-\overline v)
                            (v_\gamma-\overline v)^\top.
 \label{eq:rg-thermo-label-tie-limit}
\end{equation}
The rank is the dimension of the span of $v_\gamma-v_\eta$ over pairs
with positive limiting weights, at most
$\min(d,|\operatorname{supp}w|-1)$. A direction with positive variance
has Fisher information asymptotic to a positive multiple of $V^2$.
Zero variance implies only $o(V^2)$ information in that direction.\thermoestablishedstatus


\paragraph{Proof.}
There is a positive bulk gap from $S$ to every nonmaximizing label,
since the catalog is finite. Comparing to any one maximizing label and
using convergence of $\lambda/V$ shows that every nonmaximizing weight
is exponentially small at $\theta_0$. Thus its limit is zero. Also
$\xi_{\gamma,L}/V\to v_\gamma$ by $C^1$ bulk convergence. Dividing the
complete finite decomposition by $V^2$ makes the branch contribution
$O(1/V)$ and leaves
$\operatorname{Cov}_{W_L}(\xi_{\gamma,L}/V)$, which converges by finite
summation to \cref{eq:rg-thermo-label-tie-limit}. For any vector $z$,
\begin{equation}
 z^\top C_w(v)z
 =\sum_\gamma w_\gamma
       \big[z\cdot(v_\gamma-\overline v)\big]^2.
 \label{eq:rg-thermo-label-tie-direction}
\end{equation}
This vanishes exactly when $z\cdot v_\gamma$ is constant on the
positive-weight labels. The kernel is the orthogonal complement of
the stated span, proving the rank. A positive limiting quadratic form
gives the asserted asymptotic by multiplication by $V^2$; a zero one
gives only the stated little-oh conclusion. $\square$

\propositionheading{Residual gradients and rare labels before score squaring}
{prop:rg-thermo-label-residual-scales}
At every finite volume the label contribution has the exact form
\begin{equation}
 \frac1V\operatorname{Cov}_{W_L}(\xi_{\gamma,L})
 =\frac1{2V}\sum_{\gamma,\eta}W_{\gamma,L}W_{\eta,L}
   (\xi_{\gamma,L}-\xi_{\eta,L})
   (\xi_{\gamma,L}-\xi_{\eta,L})^\top.
 \label{eq:rg-thermo-label-pairwise-covariance}
\end{equation}
Consequently, if the weighted branch-density term converges, convergence
of the full density is equivalent to convergence of the matrix on the
right. Its weights cannot be replaced by their limits before its
score differences are squared.

Write $\lambda_{\gamma,L}=Vs_\gamma+r_{\gamma,L}$ with
$\|r_{\gamma,L}\|_{C^3(K)}/V\to0$. In the setting of
\cref{prop:rg-thermo-label-tie-scale}, suppose every label of $S$ has
positive limiting weight. Fix a direction $z$ for which $z\cdot v_\gamma$
is the same for all $\gamma\in S$, choose $\gamma_0\in S$, and assume
for some $0\leq\kappa<1$ that
\begin{equation}
 V^{-\kappa}\partial_z(r_{\gamma,L}-r_{\gamma_0,L})(\theta_0)
 \longrightarrow d_\gamma\quad(\gamma\in S),
 \qquad d_{\gamma_0}=0.
 \label{eq:rg-thermo-label-residual-assumption}
\end{equation}
Then the directional label information divided by $V^{2\kappa}$ tends
to $\operatorname{Var}_w(d_\gamma)$. Positive variance is below the
extensive scale if $\kappa<1/2$, contributes a finite nonzero density
if $\kappa=1/2$, and makes the density diverge if $\kappa>1/2$.
If all the vector slopes on $S$ agree and the corresponding relative
residual gradients divided by $\sqrt V$ tend to vectors $d_\gamma$,
then, provided $g_{\gamma,L}(\theta_0)\to g_\gamma(\theta_0)$,
\begin{equation}
 g_L(\theta_0)\longrightarrow
       \sum_{\gamma\in S}w_\gamma g_\gamma(\theta_0)
                           +\operatorname{Cov}_w(d_\gamma).
 \label{eq:rg-thermo-label-residual-density}
\end{equation}
These are fixed-parameter limits, with no conclusion about derivatives
of the limiting tensor. If a tied weight vanishes, the exact pairwise
formula still applies, but this positive-weight residual conclusion
requires separate control of that rare label.\thermoestablishedstatus


\paragraph{Proof.}
Expand the right side of \cref{eq:rg-thermo-label-pairwise-covariance}.
The two squared terms together give
\[
 V^{-1}\sum_\gamma W_\gamma\xi_\gamma\xi_\gamma^\top,
\]
and the two cross terms together give
\[
 -V^{-1}\left(\sum_\gamma W_\gamma\xi_\gamma\right)
        \left(\sum_\gamma W_\gamma\xi_\gamma\right)^\top.
\]
This proves the identity and the convergence equivalence after adding
the branch term.

On $S$, equal bulk directional slopes cancel exactly from directional
score differences at $\theta_0$. Divide those differences by
$V^\kappa$ and use \cref{eq:rg-thermo-label-residual-assumption} in the
finite pairwise sum. The limiting sum is
$\frac12\sum_{\gamma,\eta\in S}w_\gamma w_\eta(d_\gamma-d_\eta)^2
=\operatorname{Var}_w(d_\gamma)$. Pairs involving a nonmaximizing label
contribute at most $CV^{2-2\kappa}e^{-cV}$ after this normalization,
since all original score differences are $O(V)$; they vanish even when
$\kappa=0$. Comparison of $2\kappa$ with one proves the threshold
claims. Repeat the same calculation with outer products and
$\kappa=1/2$ for \cref{eq:rg-thermo-label-residual-density}; finite
summation of the branch-density limits gives its remaining term.
No derivative of any fixed-point limit has been taken. $\square$

\subsection{Explicit obstructions inside the admitted label-growth domain}
\label{sec:rg-thermo-label-counterexamples}

\propositionheading{Native realizations with divergent or nonconvergent ties}
{prop:rg-thermo-label-tie-counterexamples}
The growth condition \cref{eq:rg-thermo-label-raw-growth}, even with
identical positive high-temperature native branches, does not imply a
finite extensive Fisher limit at a tie, convergence of the label law,
or $C^2$ convergence of a uniformly convergent Fisher density. The
following families give explicit counterexamples to these assertions.\thermoestablishedstatus


\paragraph{Proof.}
Use two labels $0,1$ and identical complete native free laws at fixed
$0<b_0<1$ and $0<t_0<1/12$. The positive coupling $t_0$ is independent
of the statistical parameter $\theta$ in a fixed open interval about
zero. Every branch mass is positive on its complete carrier; its
conditional Fisher tensor in $\theta$ is zero. Its physical partition
is common to the two labels, so it cancels from their relative effective
mass in either normalization mode. Set $a_{0,L}=1$ and
$a_{1,L}=e^{f_L(\theta)}$. The binary weight and its complete Fisher
information are therefore exactly
\begin{equation}
 W_{1,L}=\frac{e^{f_L}}{1+e^{f_L}},\qquad
 I_L=q(f_L)(f'_L)^2,\qquad
 q(x)=\frac{e^x}{(1+e^x)^2}=\frac1{4\cosh^2(x/2)}.
 \label{eq:rg-thermo-label-binary-formula}
\end{equation}
Indeed $W'_1=W_1(1-W_1)f'_L$ and the two centered label scores are
$(1-W_1)f'_L$ and $-W_1f'_L$. Squaring and averaging gives this
formula. Thus no pressure derivative theorem is needed to evaluate
these examples. All assertions about ties concern the effective bulk
scores, whose common native pressure has no $\theta$ dependence.

First choose $f_L=V\theta$. Its $C^3$ norm on any compact interval is
at most $C_KV$, and $f_L/V=\theta$ exactly. At the tie $\theta=0$,
\begin{equation}
 I_L(\theta)=\frac{V^2}{4\cosh^2(V\theta/2)},\qquad
 \frac{I_L(0)}V=\frac V4.
 \label{eq:rg-thermo-label-linear-tie-example}
\end{equation}
The extensive density diverges. The pointwise $V^{-2}I_L$ limit equals
$1/4$ at zero and zero at every nonzero $\theta$, so even that pointwise
rescaling is discontinuous across the fixed-parameter tie.

Next take $f_L=0$ for even $L$ and $f_L=\log2$ for odd $L$.
Its derivatives vanish, its $C^3$ norm divided by $V$ tends to zero,
and $I_L=0$ identically. But $W_{1,L}$ alternates between $1/2$ and
$2/3$. The retained event $\{\gamma=1\}$ already fails to converge,
regardless of the changing microscopic carriers. Thus convergence of
Fisher information does not force convergence of the labeled law.

For $0<\kappa<1$ take $f_L=V^\kappa\theta$. Its $C^3$ norm divided
by $V$ is at most $C_KV^{\kappa-1}\to0$, so the two bulk scores are
identical throughout the interval. At zero both weights equal $1/2$,
and
\begin{equation}
 \frac{I_L(0)}V=\frac14V^{2\kappa-1}.
 \label{eq:rg-thermo-label-subextensive-example}
\end{equation}
This realizes all three thresholds in
\cref{prop:rg-thermo-label-residual-scales}. With
$f_L=c_L+\sqrt V\theta$, where $c_L$ alternates between $0$ and
$\log2$, the same growth bound holds while the density at zero
alternates between $q(0)=1/4$ and $q(\log2)=2/9$. Hence even bounded
extensive density need not converge.

A rare tied label is obtained with
$f_L=V\theta-a\log V$, for any fixed $a>0$. Here
$f_L/V\to\theta$ in $C^3$ and its $C^3$ norm is at most $C_KV$ for
large $V$. At zero $W_{1,L}=(1+V^a)^{-1}\to0$, but
\begin{equation}
 \frac{I_L(0)}V=\frac{V^{1+a}}{(1+V^a)^2}
                 \sim V^{1-a}.
 \label{eq:rg-thermo-label-rare-example}
\end{equation}
Thus the limiting $V^2$ coefficient is zero while the density diverges
for $a<1$, tends to one for $a=1$, and vanishes for $a>1$. Discarding
the rare label before squaring its score would miss all of these cases.

Finally put
\begin{equation}
 f_L(\theta)=A_L\sin(k_L\theta),\qquad
 A_L=V^{-1/6},\qquad k_L=V^{1/3}.
 \label{eq:rg-thermo-label-oscillatory-data}
\end{equation}
The largest bound among its derivatives through degree three is
$A_Lk_L^3=V^{5/6}$, so
$\|f_L\|_{C^3(K)}/V=O(V^{-1/6})\to0$. Positivity of the raw weights
is automatic, and $|f_L|\leq A_L$ makes both label weights converge
uniformly to $1/2$. The density
$F_L=V^{-1}q(f_L)(f'_L)^2$ obeys
$0\leq F_L\leq V^{-2/3}/4$, so it converges uniformly to zero.
Differentiate this expression twice before passing to a limit:
\begin{equation}
 F_L''=\frac1V\left[q''(f_L)(f'_L)^4
       +5q'(f_L)(f'_L)^2f_L''
       +2q(f_L)(f_L'')^2+2q(f_L)f'_Lf_L'''\right].
 \label{eq:rg-thermo-label-oscillatory-derivative}
\end{equation}
At zero, $f_L=0$, $f'_L=A_Lk_L$, $f_L''=0$, and
$f_L'''=-A_Lk_L^3$. Also $q(0)=1/4$, $q'(0)=0$, and
$q''(0)=-1/8$, obtained by differentiating
$[4\cosh^2(x/2)]^{-1}$. Consequently
\begin{equation}
 F_L''(0)=-\frac{A_L^2k_L^4}{2V}
          -\frac{A_L^4k_L^4}{8V}
          =-\frac12-\frac18V^{-1/3}\longrightarrow-\frac12.
 \label{eq:rg-thermo-label-oscillatory-failure}
\end{equation}
The uniform limit has second derivative zero. This proves failure of
$C^2$ convergence under the original $C^3$ bulk growth and convergent
finite weights. Every example has smooth positive finite data and
fixed strictly positive native coupling; none is a product-boundary
exception. $\square$

\propositionheading{A scalar boundary estimate does not control information}
{prop:rg-thermo-label-boundary-obstruction}
Uniform subvolume bounds on the log ratio of two finite laws, even
with a common bulk-cylinder limit, do not imply convergence of Fisher
densities or their parameter derivatives.\thermoestablishedstatus


\paragraph{Proof.}
Let $P_L$ be the fair product law on $V=L^2$ signs, independent of
$\theta$, and choose one sign at the outer boundary. Let $Q_{L,\theta}$
change only that sign to mean $a\sin(V\theta)$, with fixed
$0<a<1/2$. Every finite law is positive and smooth on the same open
parameter interval. Pointwise
$Q_{L,\theta}/P_L=1+a\sin(V\theta)z_{\rm boundary}$, whence
\begin{equation}
 \left|\log\frac{Q_{L,\theta}}{P_L}\right|
 \leq\max\{-\log(1-a),\log(1+a)\}.
 \label{eq:rg-thermo-label-boundary-log-ratio}
\end{equation}
Every fixed bulk cylinder eventually avoiding that boundary sign has
exactly the same law under $P_L$ and $Q_{L,\theta}$. A binary sign with
mean $m(\theta)$ has Fisher information
$(m')^2/(1-m^2)$, obtained directly by summing the squared scores of
$(1\pm m)/2$. Thus
\begin{equation}
 I(P_L)=0,\qquad
 I(Q_{L,\theta})
 =\frac{a^2V^2\cos^2(V\theta)}{1-a^2\sin^2(V\theta)},\qquad
 \frac{I(Q_{L,0})}{V}=a^2V.
 \label{eq:rg-thermo-label-boundary-information}
\end{equation}
The scalar log-ratio bound cannot transfer the density limit. This
example deliberately violates the native uniform local-response
bounds, so it refutes that inference and does not refute the conditional
native inputs. The still simpler functions
$L^{-1}\sin(L^2\theta)$ converge uniformly to zero while their first
derivatives at zero equal $L$; a scalar convergence assertion has no
implicit permission to differentiate it. $\square$

\subsection{The reciprocal-volume statistical experiment}
\label{sec:rg-thermo-label-crossover}

\theoremheading{Actual center ratios, smooth weights, and the full Fisher limit}
{thm:rg-thermo-label-crossover}
Assume \cref{eq:rg-thermo-label-effective-growth} and actual branch
limits $g_{\gamma,L}\to g_\gamma$ in $C^2$ on a neighborhood of
$\theta_0$. Let $S$ be the maximizing set of $s_\gamma(\theta_0)$,
choose $\gamma_0\in S$, and impose the additional actual center-ratio
condition
\begin{equation}
 \exp[\lambda_{\gamma,L}(\theta_0)
                 -\lambda_{\gamma_0,L}(\theta_0)]
 \longrightarrow\kappa_\gamma\in(0,\infty)
 \quad(\gamma\in S),\qquad \kappa_{\gamma_0}=1.
 \label{eq:rg-thermo-label-center-ratios}
\end{equation}
Set $v_\gamma=\nabla s_\gamma(\theta_0)$ and define the actual
finite experiment on its original labeled carrier by
\begin{equation}
 \theta_L(u)=\theta_0+u/V,\qquad
 \mathcal E_L(u)=P_{L,\theta_L(u)}.
 \label{eq:rg-thermo-label-window}
\end{equation}
On every fixed compact $u$-set, its label law converges in $C^3$ to
\begin{equation}
 \pi_\gamma(u)=
 \frac{\kappa_\gamma e^{u\cdot v_\gamma}}
      {\sum_{\eta\in S}\kappa_\eta e^{u\cdot v_\eta}}
 \quad(\gamma\in S),\qquad
 \pi_\gamma(u)=0\quad(\gamma\notin S).
 \label{eq:rg-thermo-label-limiting-weights}
\end{equation}
Its full Fisher tensor in the coordinate $u$ satisfies
\begin{equation}
 I_L^{(u)}(u)=V^{-2}I(P_{L,\theta_L(u)})
 \longrightarrow
 G(u)=\operatorname{Cov}_{\pi(u)}(v_\gamma)
 =D_u^2\log\sum_{\gamma\in S}\kappa_\gamma e^{u\cdot v_\gamma}
 \quad\hbox{in }C^2.
 \label{eq:rg-thermo-label-crossover-fisher}
\end{equation}
The center-ratio condition is additional finite-volume data. Equality of
bulk pressure or effective bulk scores supplies no such condition.\thermoestablishedstatus


\paragraph{Proof.}
A fixed compact $u$-set has bounded radius, so its image and the line
segments to $\theta_0$ lie in the chosen parameter neighborhood for
large $L$. Define
$\Delta_{\gamma,L}(u)=\lambda_{\gamma,L}(\theta_0+u/V)
-\lambda_{\gamma,L}(\theta_0)$. The fundamental theorem of calculus
on each segment gives
\begin{equation}
 \Delta_{\gamma,L}(u)
 =\int_0^1 V^{-1}\nabla\lambda_{\gamma,L}(\theta_0+\tau u/V)
                                  \cdot u d\tau.
 \label{eq:rg-thermo-label-window-increment}
\end{equation}
The integrand tends uniformly to $v_\gamma\cdot u$, by the compact
$C^1$ convergence in \cref{eq:rg-thermo-label-effective-growth} and
continuity of $\nabla s_\gamma$. Direct differentiation in $u$ yields
\begin{equation}
 D_u\Delta_{\gamma,L}\to v_\gamma,\qquad
 \|D_u^2\Delta_{\gamma,L}\|=O(V^{-1}),\qquad
 \|D_u^3\Delta_{\gamma,L}\|=O(V^{-2}).
 \label{eq:rg-thermo-label-window-derivatives}
\end{equation}
Thus $\Delta_{\gamma,L}\to u\cdot v_\gamma$ in $C^3$. For tied labels,
subtract the reference increment and add the log of the actual center
ratio. The resulting relative log masses converge in $C^3$ to
$\log\kappa_\gamma+u\cdot(v_\gamma-v_{\gamma_0})$.
Exponentiation and the finite softmax are smooth on the resulting
compact ranges; their denominator is bounded below by the reference
mass, which equals one in relative coordinates.

For a nonmaximizing label the relative center log mass is at most
$-cV$ for some $c>0$ and large $L$. Its window increment is $O(1)$
and its $u$-derivatives through degree three are bounded. Differentiating
its relative exponential through degree three therefore bounds each
term by $Ce^{-cV}$. These masses and all those derivatives vanish.
Dividing the complete finite vector of relative masses by its sum
proves \cref{eq:rg-thermo-label-limiting-weights} in $C^3$, including
the zero limiting coordinates outside $S$.

Define the normalized effective score
$\zeta_{\gamma,L}(u)=V^{-1}\xi_{\gamma,L}(\theta_L(u))$.
Equation \cref{eq:rg-thermo-label-window-derivatives} implies
$\zeta_{\gamma,L}\to v_\gamma$ in $C^2$, for all labels.
The affine parameter change in \cref{eq:rg-thermo-label-window}
has derivative $V^{-1}\operatorname{Id}$, so the full finite Fisher
identity becomes
\begin{equation}
 I_L^{(u)}(u)=\operatorname{Cov}_{W_L(\theta_L(u))}
                         (\zeta_{\gamma,L}(u))
       +\frac1V\sum_\gamma W_{\gamma,L}(\theta_L(u))
                                g_{\gamma,L}(\theta_L(u)).
 \label{eq:rg-thermo-label-window-decomposition}
\end{equation}
The covariance converges in $C^2$ by finite sums and products. This
form never divides by a weight tending to zero. The branch densities
and their first two $u$-derivatives are bounded, since a $u$ derivative
is $V^{-1}$ times a $\theta$ derivative. The weights have bounded $C^2$
norms by the preceding argument. Every term in the ordinary
second-order Leibniz sum for the last term in
\cref{eq:rg-thermo-label-window-decomposition} is therefore $O(V^{-1})$.
The conditional-state information vanishes in the limit but is not
identically zero in general. Finally, differentiating the finite
log-sum-exponential gives gradient $\sum\pi_\gamma v_\gamma$ and
Hessian $\operatorname{Cov}_\pi(v)$, completing the proof. $\square$

\propositionheading{Two parameter-independent comparison kernels}
{prop:rg-thermo-label-comparison-kernels}
Under \cref{thm:rg-thermo-label-crossover}, let $\Pi(u)$ denote the
finite-label experiment on the full catalog with masses
$\pi_\gamma(u)$, including its structural zeros. There are kernels in
both directions between $\mathcal E_L(u)$ and $\Pi(u)$ that are
independent of $u$ and have uniformly vanishing total-variation errors
on every compact $u$-set. The replacement of a conditional branch by
its actual center law has error $O(V^{-1/2})$; the additional label
error has only the convergence rate furnished by
\cref{eq:rg-thermo-label-center-ratios}.\thermoestablishedstatus


\paragraph{Proof.}
Let $K_L$ extract the label and let $R_L$ sample the corresponding
actual finite center branch:
\begin{equation}
 K_L(\delta\mid\gamma,x)=\boldsymbol1_{\{\delta=\gamma\}},\qquad
 R_L(\delta,x\mid\gamma)=\boldsymbol1_{\{\delta=\gamma\}}
                                      p_{\gamma,L,\theta_0}(x).
 \label{eq:rg-thermo-label-comparison-kernels}
\end{equation}
Both are normalized, preserve labels, and may depend on $L$ but not on
$u$. The center branch is defined for every finite positive label,
even when its limiting label mass is zero. Put
$\varepsilon_L=\sup_u\frac12\sum_\gamma
 |W_{\gamma,L}(\theta_L(u))-\pi_\gamma(u)|$ over the chosen compact.
The first comparison has error at most $\varepsilon_L\to0$.

For the reverse comparison, each branch has
$I(p_{\gamma,L,\theta})\preceq CV\operatorname{Id}$ along the
segments in question, by its bounded density. On the fixed active
support, direct differentiation gives
\begin{equation}
 \left\|\frac{d}{d\tau}
        \sqrt{p_{\gamma,L,\theta_0+\tau u/V}}\right\|_2^2
 =\frac1{4V^2}u^\top
       I(p_{\gamma,L,\theta_0+\tau u/V})u
 \leq\frac{C|u|^2}{4V}.
 \label{eq:rg-thermo-label-square-root-bound}
\end{equation}
The square root vector is differentiable because the support is fixed
and positive; zero coordinates can simply be left identically zero.
Integrating its derivative and using the triangle inequality gives
square-root Euclidean distance at most $C'|u|/\sqrt V$ between the
branch at $\theta_L(u)$ and its center. For any probability vectors
$p,q$, Cauchy-Schwarz gives
\begin{equation}
 \|p-q\|_{\rm TV}
 =\frac12\sum_x|\sqrt{p(x)}-\sqrt{q(x)}|
                      (\sqrt{p(x)}+\sqrt{q(x)})
 \leq\|\sqrt p-\sqrt q\|_2,
 \label{eq:rg-thermo-label-tv-square-root}
\end{equation}
since $\|\sqrt p+\sqrt q\|_2\leq2$. On the disjoint labeled carrier,
split $W_\gamma p_{\gamma,\theta_L}-\pi_\gamma p_{\gamma,\theta_0}$
into $W_\gamma(p_{\gamma,\theta_L}-p_{\gamma,\theta_0})$
and $(W_\gamma-\pi_\gamma)p_{\gamma,\theta_0}$, then sum absolute
values. This proves the explicit bound
\begin{equation}
 \sup_u\|\mathcal E_L(u)-\Pi(u)R_L\|_{\rm TV}
 \leq C_KV^{-1/2}+\varepsilon_L.
 \label{eq:rg-thermo-label-experiment-error}
\end{equation}
This constructs comparison of actual experiments without an infinite
microscopic partition sum or a common limiting state carrier.
$\square$

\propositionheading{The actual Euclidean null-slope quotient}
{prop:rg-thermo-label-quotient}
Let
\begin{equation}
 E=\operatorname{span}\{v_\gamma-v_{\gamma_0}:\gamma\in S\},
 \qquad r=\dim E.
 \label{eq:rg-thermo-label-slope-space}
\end{equation}
For every $u$, $\ker G(u)=E^\perp$, and
$\pi(u)=\pi(u')$ holds exactly when $u-u'\in E^\perp$.
The limiting family therefore descends to the Euclidean quotient
$\R^d/E^\perp$, identified with $E$. On $E$ its probability map is
injective and, if $r>0$, its Fisher tensor is positive definite with a
positive minimum on every compact chart. The restrictions of the
finite $u$-metrics to this fixed subspace converge in $C^2$ and have
the inverse, connection, and curvature limits of
\cref{prop:rg-thermo-label-regular-geometry}.\thermoestablishedstatus


\paragraph{Proof.}
All limiting weights on $S$ are positive at every finite $u$.
The directional quadratic form of $G$ is their variance of
$z\cdot v_\gamma$, so it vanishes exactly when $z$ is orthogonal to
all slope differences. This proves the kernel statement. The
probability ratios satisfy the exact formula
\begin{equation}
 \log\frac{\pi_\gamma(u)}{\pi_{\gamma_0}(u)}
 =\log\kappa_\gamma+u\cdot(v_\gamma-v_{\gamma_0}).
 \label{eq:rg-thermo-label-quotient-ratios}
\end{equation}
Equal laws therefore imply orthogonality of $u-u'$ to every generating
slope difference, and that orthogonality conversely makes every ratio,
hence every probability, equal. Orthogonal projection onto $E$ gives
the claimed quotient identification. Choosing $r$ independent slope
differences gives a nonsingular linear map on $E$; their log ratios
recover the coordinates on $E$, which also proves injectivity and a
smooth inverse onto the image of the probability map.

Choose orthonormal coordinates on $E$. The restriction of $G$ is
positive definite, since its only null vector in $E$ is zero. Its
continuous quadratic form attains a positive minimum on each compact
chart times the unit sphere in $E$. Constant linear restriction
preserves the $C^2$ convergence in
\cref{eq:rg-thermo-label-crossover-fisher}, and
\cref{prop:rg-thermo-label-regular-geometry} applies. When $r=0$ the
limiting experiment is constant and the quotient is a point; no
positive-dimensional inverse is asserted. The finite metrics need not
descend along $E^\perp$: this statement restricts them to a chosen
fixed copy of $E$. Slower information scales in the other directions
are not excluded. $\square$

\propositionheading{A nonempty native realization of the limiting experiment}
{prop:rg-thermo-label-native-witness}
Every finite list of slopes $v_\gamma\in\R^d$ and positive constants
$\kappa_\gamma$, with $\kappa_{\gamma_0}=1$, has a realization on
identical native branches with strictly positive high-temperature
plaquette coupling for which the crossover weights, Fisher tensor,
and both comparison kernels above are exact at every finite volume.\thermoestablishedstatus


\paragraph{Proof.}
Choose common fixed $0<b_0<1$, $0<t_0<1/12$ and the complete native
free branch for every label. Set
$a_{\gamma,L}(\theta)=\kappa_\gamma e^{Vv_\gamma\cdot\theta}$.
The weights are positive and smooth, their log $C^3$ norms on compact
sets are $O(V)$, and their bulk logarithms converge in $C^3$ to
$v_\gamma\cdot\theta$. The identical physical partitions cancel in
all label ratios, both globally and hierarchically. At $\theta_0=0$
the actual center ratios are $\kappa_\gamma$ exactly, and at $u/V$
the normalized label weights are precisely
\cref{eq:rg-thermo-label-limiting-weights}. Conditional branches are
independent of $\theta$, so their Fisher term is zero and the
reconstruction kernel uses exactly their unchanged law. Both
comparison errors are zero. With two labels in one coordinate,
$v_0=0$, $v_1=1$, and $\kappa_0=\kappa_1=1$, this gives
\begin{equation}
 I_L^{(u)}=G(u)=\frac1{4\cosh^2(u/2)}>0.
 \label{eq:rg-thermo-label-binary-native-window}
\end{equation}
Its one-dimensional intrinsic curvature is zero, since the two
antisymmetrized derivative indices in the Riemann formula must
coincide. The branch coupling is strictly positive; only its dependence
on the crossover coordinate is absent. $\square$

\theoremheading{A transverse two-label moving center}
{thm:rg-thermo-label-moving-center}
Assume the effective growth and actual branch $C^2$ inputs of
\cref{thm:rg-thermo-label-crossover}. Suppose exactly two labels
$0,1$ maximize at $\theta_0$, all others have a strict gap there, and
for a fixed unit vector $n$,
$a:=\partial_n(s_1-s_0)(\theta_0)\ne0$. Reverse $n$ if necessary so
$a>0$. On a sufficiently short fixed segment through $\theta_0$ there
is, for all large $L$, exactly one center
$\theta_{c,L}=\theta_0+t_Ln$ satisfying
$\lambda_{1,L}(\theta_{c,L})=\lambda_{0,L}(\theta_{c,L})$.
It obeys
\begin{equation}
 |t_L|\leq\frac{4|r_{1,L}(\theta_0)-r_{0,L}(\theta_0)|}{Va}
 \longrightarrow0,
 \qquad r_{\gamma,L}=\lambda_{\gamma,L}-Vs_\gamma.
 \label{eq:rg-thermo-label-moving-center-bound}
\end{equation}
The experiment $P_{L,\theta_{c,L}+u/V}$ has the $C^3$ weight and $C^2$
full Fisher limits of \cref{thm:rg-thermo-label-crossover}, with
$S=\{0,1\}$, $\kappa_0=\kappa_1=1$, and slopes
$v_\gamma=\nabla s_\gamma(\theta_0)$. Its comparison kernel samples
$p_{\gamma,L,\theta_{c,L}}$, independently of $u$.\thermoestablishedstatus


\paragraph{Proof.}
Let $f=s_1-s_0$. Choose $\rho>0$ such that the closed segment
$\theta_0+[-\rho,\rho]n$ is inside the parameter neighborhood,
$\partial_nf\geq a/2$ on it, and both labels $0,1$ remain uniformly
above every other bulk score there. Such a choice exists by continuity,
the nonzero transverse derivative, and the finite set of strict gaps
at $\theta_0$. Since $f(\theta_0)=0$, integrating its derivative gives
$f(\theta_0+\rho n)\geq a\rho/2$ and
$f(\theta_0-\rho n)\leq-a\rho/2$.

For $F_L=\lambda_{1,L}-\lambda_{0,L}$, $C^1$ bulk convergence gives
$\partial_nF_L\geq Va/4$ throughout this segment for large $L$.
Uniform bulk convergence also gives opposite signs of $F_L$ at its
endpoints. Continuity supplies a zero and strict monotonicity makes
it unique. Since
$F_L(\theta_0)=r_{1,L}(\theta_0)-r_{0,L}(\theta_0)=o(V)$, integration
of the derivative lower bound between zero and $t_L$ gives
\cref{eq:rg-thermo-label-moving-center-bound}. In particular the centers
converge to $\theta_0$ and eventually lie in its interior neighborhood.

At the center the relative mass of label one to zero is exactly one.
Other center relative masses are exponentially small by the uniform
bulk gaps. For compact $u$-sets, apply the segment identity
\cref{eq:rg-thermo-label-window-increment} with base point
$\theta_{c,L}$ instead of $\theta_0$. Compact $C^3$ growth and
$\theta_{c,L}\to\theta_0$ give the same first-derivative limits and
second/third derivative bounds in
\cref{eq:rg-thermo-label-window-derivatives}. The actual centered
relative masses therefore converge in $C^3$ to
$\exp[u\cdot(v_1-v_0)]$, with exponentially vanishing complementary
masses. The covariance proof and the $O(V^{-1})$ branch bound in
\cref{eq:rg-thermo-label-window-decomposition} give $C^2$ full Fisher
convergence. Finally use the two kernels in
\cref{eq:rg-thermo-label-comparison-kernels} with the actual moving
center branch. The square-root argument is uniform on the neighborhood
and again gives the bound in
\cref{eq:rg-thermo-label-experiment-error}. $\square$

The displacement bound permits $|t_L|$ much larger than $V^{-1/2}$.
Sampling the fixed $\theta_0$ branch instead would require control of
that additional displacement: the same information estimate bounds
its total-variation cost only by $C\sqrt V|t_L|$, which need not vanish.
No such replacement is used. More than two tied labels would require
solving several relative-mass equations simultaneously; the single
transverse argument does not assert that solvability. The reciprocal
coordinate $u$ defines a statistical experiment at a declared volume
scale. It is neither physical time nor a universal smooth continuation
of the extensive tensor across arbitrary ties.

% Intended for the existing Theory semantic headings and shared macros.
% No preamble, chapter heading, status assignment, or manuscript input is added.

\section{Fixed-word scales, finite jets, and prescribed actions}
\label{sec:rg-thermo-accounting-proof}

The exact retained experiment fixes the meaning of every scale comparison.
Throughout this section, $\Theta\subset\R^d$ is open and the geometry
catalog $\mathfrak G$ is finite and nonempty. At each finite volume,
$P_\theta(\gamma,x)=W_\gamma(\theta)p_{\gamma,\theta}(x)$ is a
normalized law on the disjoint union of finite branch carriers. Its
positive support is fixed on the parameter neighborhood under discussion,
and its masses are $C^3$. Conditional objects are defined only on
occupied fibers. A channel is normalized, independent of $\theta$, and
preserves $\gamma$. The source may be a prior or a positive selected
posterior. For a random observation joint, the observation is retained
as part of both endpoint variables. These conventions are those of
\Cref{def:rg-thermo-finite-channel}; they exclude replacing the induced
observation model by a new local sensor.\thermoadoptedhypothesis


\propositionheading{Finite cell-area accounting}
{prop:rg-thermo-account-area}
Fix a finite block word with total linear factor $B\geq1$. Its fine and
final coarse cell counts are $M=B^2V$ and $V$, respectively. Let $I_f$
and $I_c$ be the full Fisher tensors of one finite experiment and its
exact image, evaluated in the same parameter coordinates, and put
$D=I_f-I_c$. Then
\begin{align}
 B^2\frac{I_f}{M}&=\frac{I_c}{V}+\frac{D}{V},
 \label{eq:rg-thermo-account-extensive}\\
 B^4\frac{I_f}{M^2}&=\frac{I_c}{V^2}+\frac{D}{V^2}.
 \label{eq:rg-thermo-account-quadratic}
\end{align}
The loss is positive semidefinite, and more precisely
\begin{equation}
 \begin{aligned}
 D&=\sum_\gamma W_\gamma
    \bigl[I(p_\gamma)-I(p_\gamma C_\gamma)\bigr],\\
 0&\preceq D\preceq\sum_\gamma W_\gamma I(p_\gamma).
 \end{aligned}
 \label{eq:rg-thermo-account-branch-bound}
\end{equation}
All identities have their ordinary parameter derivatives through order
two; the sign assertions concern the tensors themselves.\thermoestablishedstatus


\paragraph{Proof.}
The conditional-score identity of
\Cref{prop:rg-thermo-finite-fisher-loss} gives $D\succeq0$ and its
weighted branch representation. Each pushed branch Fisher tensor is
positive semidefinite because it is a score Gram matrix. Subtracting
these tensors from the weighted fine branch sum proves the upper bound
in \cref{eq:rg-thermo-account-branch-bound}. Division of the exact
identity $I_f=I_c+D$ by $V$ and by $V^2$, followed by $M=B^2V$,
gives \cref{eq:rg-thermo-account-extensive,eq:rg-thermo-account-quadratic}.
Volumes and $B$ are fixed integers when parameter derivatives are taken,
so the same equalities hold after each available derivative. No sign
for those derivative matrices follows by differentiating an order
inequality. Two edge signs per cell change the native statistic counts,
not the relation between the two cell areas. $\square$

If $I_f/M\to g_f$ and $I_c/V\to g_c$ in $C^2(K)$ on a compact
parameter set $K$, then the first identity gives
$D/V\to B^2g_f-g_c$ in $C^2(K)$, and its limiting value is positive
semidefinite. This conclusion requires both density limits as inputs.
The fixed word precedes every volume limit in this section.

For the native family, these branch inputs are supplied by
\Cref{prop:rg-thermo-native-ctwo-metric,thm:rg-thermo-native-free-push-ctwo}.
Precisely, let each $\chi_\gamma:\Theta\to\R^2$ be $C^4$ on a
neighborhood of $K$, with image in a common compact source box of
\Cref{hyp:rg-thermo-native-source-box}, and take
$p^f_{\gamma,M}=\nu^{\mathrm f}_{BL,\chi_\gamma}$ and
$p^c_{\gamma,L}=(C_B)_\#\nu^{\mathrm f}_{BL,\chi_\gamma}$.
Here $M=(BL)^2$ and the word is the exact parity word of
\Cref{def:rg-thermo-native-families}. The source-score chain rule and
those two results give
\begin{equation}
 \begin{aligned}
 \frac{I(p^f_{\gamma,M})}{M}
 &\longrightarrow D\chi_\gamma^\top(H\circ\chi_\gamma)D\chi_\gamma,\\
 \frac{I(p^c_{\gamma,L})}{V}
 &\longrightarrow
 D\chi_\gamma^\top(\mathcal I_B\circ\chi_\gamma)D\chi_\gamma
 \end{aligned}
 \quad\text{in }C^2(K).
 \label{eq:rg-thermo-account-native-density-inputs}
\end{equation}
The coarse expression is a pulled-back exact-push density. It is not
identified with the Hessian of a native coarse pressure. Compact
$C^2$ convergence bounds the fine densities uniformly; finitely many
remaining small sizes are bounded by finite smoothness and compactness.
Taking the maximum over the finite catalog supplies
\cref{eq:rg-thermo-account-extensive-branch-input}.

For the full labeled tensors, also impose the effective $C^3$ growth
and strict-gap hypotheses of \Cref{thm:rg-thermo-label-strict}, with
fine area $M$ in place of its volume $V$. The same representative
$\lambda^f_{\gamma,M}$ can be used after the label-preserving push.
Consequently $\lambda^c_{\gamma,L}/V=B^2\lambda^f_{\gamma,M}/M$
converges in $C^3$ to $B^2s_\gamma$, and its $C^3$ norm is
$O(V)$ for fixed $B$. The winner is unchanged and the limiting gap is
multiplied by $B^2$. Applying the strict-phase result to each endpoint
therefore supplies $g_f$ and $g_c$ above as the corresponding winning
branch limits in \cref{eq:rg-thermo-account-native-density-inputs}.
\Cref{prop:rg-thermo-label-representatives} supplies the effective
growth from the admitted raw weights in each normalization mode; it
does not supply a strict gap. For other channels or branch families,
the two density limits remain explicit inputs: a fine limit alone does
not produce a coarse one or a canonical free coarse ensemble.\thermoestablishedstatus

\propositionheading{Preservation of a proved superextensive scale}
{prop:rg-thermo-account-superextensive}
Suppose that on a compact $K\subset\Theta$ there is a constant $C_K$
independent of $M$ and $\gamma$ such that
\begin{equation}
 I(p_{\gamma,M,\theta})\preceq C_K M\operatorname{Id}
 \quad\text{for all }\theta\in K.
 \label{eq:rg-thermo-account-extensive-branch-input}
\end{equation}
Let $a_M>0$ be a parameter-independent information scale with
$M/a_M\to0$. If $I_f/a_M\to H$ pointwise on $K$, or uniformly
there, then $I_c/a_M\to H$ in the same sense. In particular, for a
fixed exponent $p>1$, a proved limit $I_f/M^p\to H$ implies
\begin{equation}
 \frac{I_c}{V^p}\longrightarrow B^{2p}H,
 \qquad \frac{D}{V^p}\longrightarrow0.
 \label{eq:rg-thermo-account-power-limit}
\end{equation}
No convergence of parameter derivatives is asserted from
\cref{eq:rg-thermo-account-extensive-branch-input} alone.\thermoestablishedstatus


\paragraph{Proof.}
The branch bound and $\sum_\gamma W_\gamma=1$ imply
$0\preceq D\preceq C_KM\operatorname{Id}$. Consequently
$\|I_f/a_M-I_c/a_M\|\leq C_KM/a_M\to0$ in operator norm,
uniformly on $K$. This proves the first assertion with the common
fine-volume denominator $a_M$. For $a_M=M^p$, use
$M^p=B^{2p}V^p$ to obtain
\cref{eq:rg-thermo-account-power-limit}. No relation between $a_M$
and $a_V$ is implied for an arbitrary sequence $a$; the power-law
multiplier follows from the specified power itself. $\square$

For example, the finite score argument in
\Cref{prop:rg-thermo-label-tie-scale} applies with fine area $M$.
Suppose at a fixed parameter the normalized label weights
converge to $w$, the effective label log masses have
$M^{-1}\nabla\lambda_{\gamma,M}\to v_\gamma$, and
\cref{eq:rg-thermo-account-extensive-branch-input} holds. The complete
score decomposition then gives
$I_f/M^2\to\operatorname{Cov}_w(v_\gamma)$. The coarse tensor per
squared coarse area tends to
$B^4\operatorname{Cov}_w(v_\gamma)$. Likewise any separately proved
scale $M^{2\kappa}$ with $\kappa>1/2$ is preserved, with multiplier
$B^{4\kappa}$ when measured per $V^{2\kappa}$. The label marginal is
already preserved exactly at every finite volume; the branch bound
shows that the discarded state information cannot contribute at these
larger scales. Equality of bulk log-mass densities alone supplies neither the
normalized weight limit nor any such information limit.

\propositionheading{Finite-word score and VFE telescopes}
{prop:rg-thermo-account-word}
Let $C_1,\ldots,C_m$ be a fixed finite sequence of channels of the
stated kind. Put $P_j=P_0C_1\cdots C_j$, let $S_j$ be its normalized
score, and write $I_j=I(P_j)$. Then
\begin{equation}
 I_0-I_m=\sum_{j=1}^m D_j,
 \qquad D_j=I_{j-1}-I_j\succeq0.
 \label{eq:rg-thermo-account-score-telescope}
\end{equation}
For a fixed selected observation with evidence $E_o>0$, let $\Pi_0$
be the selected posterior and let $Q_0\ll\Pi_0$ be a recognition law.
Push both by this same word and set
$\Fenergy_j=-\log E_o+\KL(Q_j\Vert\Pi_j)$. If $K_j^{\mathrm{KL}}$
is the expected recognition-to-posterior conditional KL on the reverse
fiber of step $j$, then
\begin{equation}
 \Fenergy_0-\Fenergy_m=\sum_{j=1}^m K_j^{\mathrm{KL}},
 \qquad K_j^{\mathrm{KL}}\geq0.
 \label{eq:rg-thermo-account-vfe-telescope}
\end{equation}
For total area factor $B^2$, this also gives
\begin{equation}
 B^2\frac{\Fenergy_0}{M}
 =\frac{\Fenergy_m}{V}+\frac{1}{V}\sum_{j=1}^mK_j^{\mathrm{KL}}.
 \label{eq:rg-thermo-account-vfe-area}
\end{equation}\thermoestablishedstatus


\paragraph{Proof.}
Use the joint path law
$P_0(dx_0)C_1(x_0,dx_1)\cdots C_m(x_{m-1},dx_m)$, including
all retained labels and observations in each $x_j$. Finite differentiation
shows $S_j(X_j)=\E[S_0(X_0)\mid X_j]$ and
$\E[S_{j-1}(X_{j-1})\mid X_j]=S_j(X_j)$. Put
$\varepsilon_j=S_{j-1}(X_{j-1})-S_j(X_j)$. The Markov property
gives zero conditional mean of $\varepsilon_j$ given
$X_j,\ldots,X_m$. For $j<k$, $\varepsilon_k$ is measurable with
respect to these later variables, so
$\E[\varepsilon_j\varepsilon_k^\top]=0$. Thus the telescoping
score residual $S_0-S_m=\sum_j\varepsilon_j$ has second moment
$\sum_j\E[\varepsilon_j\varepsilon_j^\top]$, and each summand is
$D_j$ by the one-step conditional-score identity. This proves
\cref{eq:rg-thermo-account-score-telescope} on one augmented path
space; unrelated output spaces need not have nested sigma-algebras.

For VFE, finite domination persists under each channel, so all displayed
KL values are finite. Apply
\Cref{prop:rg-thermo-finite-vfe-loss} at each step to obtain
$\Fenergy_{j-1}-\Fenergy_j=K_j^{\mathrm{KL}}$. The evidence is
unchanged because every step pushes the complete observation experiment,
as in \Cref{prop:rg-thermo-finite-observation-push}. Summing cancels
the intermediate objectives and proves
\cref{eq:rg-thermo-account-vfe-telescope}; division by $V$ gives
\cref{eq:rg-thermo-account-vfe-area}. A new likelihood at an
intermediate stage would destroy this common-evidence argument.
The formula compares the given recognition law with its actual images;
it asserts neither a stochastic variational-infimum identity nor a
thermodynamic VFE limit. $\square$

\propositionheading{Fine and coarse crossover coordinates}
{prop:rg-thermo-account-window-coordinates}
Let $\theta_{c,L}$ be a declared center, independent of the local
coordinate, and compare
\begin{equation}
 u=M(\theta-\theta_{c,L}),\qquad
 v=V(\theta-\theta_{c,L}),\qquad u=B^2v.
 \label{eq:rg-thermo-account-window-chart}
\end{equation}
In the common $u$ coordinate, the fine and coarse tensors are
$M^{-2}I_f(\theta_{c,L}+u/M)$ and
$M^{-2}I_c(\theta_{c,L}+u/M)$. Under the uniform branch bound
on a compact neighborhood of these windows, their difference tends to
zero uniformly on compact $u$ sets. If their common limit is $H(u)$,
the coarse tensor in its own reciprocal-coarse-volume coordinate is
\begin{equation}
 I_c^{(v)}(v)=B^4 I_c^{(u)}(B^2v)
 \longrightarrow B^4H(B^2v).
 \label{eq:rg-thermo-account-window-pullback}
\end{equation}\thermoestablishedstatus


\paragraph{Proof.}
The differential of $\theta_{c,L}+u/M$ with respect to $u$ is
$M^{-1}\operatorname{Id}$, so the two factors in the Fisher pullback
give $M^{-2}$. Their loss is bounded by $C_K/M$ using
\cref{eq:rg-thermo-account-branch-bound}. This proves uniform value
convergence of the difference wherever the windows stay in the compact
parameter neighborhood. The differential of $u=B^2v$ is
$B^2\operatorname{Id}$; pulling back once more gives both the factor
$B^4$ and the argument $B^2v$ in
\cref{eq:rg-thermo-account-window-pullback}. The retained label law
at the same $u$ is identical before and after the channel at every finite
volume, not merely in its limit. $\square$

A derivative version requires derivative inputs. Write
$g_{\gamma,M}^f=I(p^f_{\gamma,M})/M$ and
$g_{\gamma,M}^c=I(p^c_{\gamma,L})/M$, using the fine denominator
for both. Suppose these tensors have uniformly bounded $C^2$ norms in
the original $\theta$ chart on a compact neighborhood of the windows,
and the reparameterized label weights have uniformly bounded $C^2$
norms on compact $u$ sets. The within-branch contribution to the loss
in the $u$ chart is then
\begin{equation}
 \frac1M\sum_\gamma W_{\gamma,M}(\theta_{c,L}+u/M)
 (g_{\gamma,M}^f-g_{\gamma,M}^c)(\theta_{c,L}+u/M).
 \label{eq:rg-thermo-account-window-loss}
\end{equation}
Twice applying the product and chain rules bounds its $C^2$ norm by
$O(M^{-1})$: every derivative of a composed branch density contributes
an additional $M^{-1}$, and the label-weight derivatives are bounded.
Under these inputs, $C^2$ convergence of the fine window tensor
transfers to the coarse one, and
\cref{eq:rg-thermo-account-window-pullback} holds in $C^2$.
The original value bound alone does not justify this upgrade.

For the native branches of
\cref{eq:rg-thermo-account-native-density-inputs}, the needed branch
bounds follow from those actual $C^2$ limits. The coarse density there
is divided by $V$; the density here is its $B^{-2}$ multiple.
For a fixed center $\theta_0$, assume the effective growth and the
actual positive center ratios in
\Cref{thm:rg-thermo-label-crossover}, using the fine area $M$.
That theorem gives $C^3$ convergence of the reparameterized weights
and the fine $C^2$ window limit
$H(u)=\operatorname{Cov}_{\pi(u)}(v_\gamma)$, where
$v_\gamma=\nabla s_\gamma(\theta_0)$ and $s_\gamma$ is normalized
per fine cell. The preceding loss bound therefore gives this same
$H(u)$ for the exact push in the common $u$ chart. In the coarse
$v$ chart the limit is $B^4H(B^2v)$; equivalently the coarse bulk
slopes are $B^2v_\gamma$. Neither the factor nor the argument may be
removed by identifying the two reciprocal-volume coordinates.

For a moving center, \Cref{thm:rg-thermo-label-moving-center} supplies
the weight and fine-tensor inputs when exactly two labels maximize at
$\theta_0$, the other labels have a strict gap there, and their bulk
log-mass density difference has a nonzero derivative along the declared transverse
direction. Apply it on the fine-square subsequence of side $BL$, and
use its actual equal-log-mass center as $\theta_{c,L}$. The weights are
unchanged by the push, so this is the same center at both endpoints.
The theorem gives unit relative center ratios and the required $C^3$
weight and $C^2$ Fisher limits; the native bounds again transfer the
limit. An arbitrary declared center, a nontransverse tie, or a tie of
more than two labels does not acquire these center properties from this
argument. For fixed centers the positive ratio condition remains an
additional premise, since the bulk pressure limit alone does not imply it.\thermoestablishedstatus

\subsection{Ordinary finite probability jets}
\label{sec:rg-thermo-account-jets}

For a positive integer $d$, write
$\mathcal O_m=\{\alpha\in\mathbb N^d:|\alpha|\leq m\}$,
$\alpha!=\prod_i\alpha_i!$, and
$\binom\alpha\beta=\prod_i\binom{\alpha_i}{\beta_i}$ for
$\beta\leq\alpha$ componentwise. Every jet below stores ordinary
derivatives at one point $\theta_0$, not Taylor coefficients. Thus a
coefficient in a Taylor polynomial is divided by $\alpha!$.\thermodeclareddefinition


\propositionheading{Local realization on fixed support}
{prop:rg-thermo-account-jet-realization}
Let $\mathsf X$ be a finite nonempty carrier and prescribe real rows
$\pi_\alpha(x)$ for $\alpha\in\mathcal O_3$. Assume
\begin{equation}
 \begin{aligned}
 \pi_0(x)&\geq0,\qquad \sum_x\pi_0(x)=1,\\
 \sum_x\pi_\alpha(x)&=0\quad(0<|\alpha|\leq3),
 \end{aligned}
 \label{eq:rg-thermo-account-jet-admission}
\end{equation}
and require $\pi_\alpha(x)=0$ for every supplied $\alpha$ whenever
$\pi_0(x)=0$. These rows are the ordinary third-order jet of a smooth
normalized family with exactly the same positive support on a
neighborhood of $\theta_0$.\thermoestablishedstatus


\paragraph{Proof.}
Set
\begin{equation}
 \pi_{\theta_0+z}(x)
 =\sum_{\alpha\in\mathcal O_3}
              \frac{\pi_\alpha(x)}{\alpha!}z^\alpha.
 \label{eq:rg-thermo-account-jet-polynomial}
\end{equation}
Summing the finite polynomial over $x$ and using
\cref{eq:rg-thermo-account-jet-admission} gives one identically. Each
atom with $\pi_0(x)>0$ remains positive on some neighborhood of
$z=0$ by continuity. There are finitely many such atoms, so the
intersection of these neighborhoods is a common open neighborhood.
Every other column is the zero polynomial. The family is therefore a
probability law with the required fixed support there, and differentiation
at zero recovers each supplied ordinary row. This construction gives a
local realization; it does not identify the polynomial with a particular
native law that might have produced the same derivatives. $\square$

The support condition is substantive. A nonnegative family may have
$\pi_{\theta_0}(x)=0$ and a nonzero second derivative at that atom
while creating positive mass nearby. Such a family is outside this
fixed-support realization and score interface. A zero atom retained only
as part of the declared carrier has no score at $\theta_0$ and contributes
zero to all score sums here.

\lemmaheading{Product, reciprocal, and normalized mass recurrences}
{lem:rg-thermo-account-jet-recurrences}
For scalar jets $a,b$ through degree $m$, their product has rows
\begin{equation}
 (ab)_\alpha=\sum_{\beta\leq\alpha}
       \binom\alpha\beta a_\beta b_{\alpha-\beta}.
 \label{eq:rg-thermo-account-jet-product}
\end{equation}
If $a_0>0$, the reciprocal rows are uniquely determined by
\begin{equation}
 r_0=a_0^{-1},\qquad
 r_\alpha=-a_0^{-1}
     \sum_{0<\beta\leq\alpha}
       \binom\alpha\beta a_\beta r_{\alpha-\beta}
       \quad(0<|\alpha|\leq m).
 \label{eq:rg-thermo-account-jet-reciprocal}
\end{equation}
For unnormalized mass rows $w_\alpha(x)$ on a fixed finite support,
let $Z_\alpha=\sum_xw_\alpha(x)$ and assume $Z_0>0$ and positive
base masses on that support. If $r$ is the reciprocal jet of $Z$, then
\begin{equation}
 \pi_\alpha(x)=\sum_{\beta\leq\alpha}
        \binom\alpha\beta w_\beta(x)r_{\alpha-\beta}
 \label{eq:rg-thermo-account-jet-normalization}
\end{equation}
is the normalized probability jet, provided structural zero columns have
all their supplied derivatives zero.\thermoestablishedstatus


\paragraph{Proof.}
The multivariate Leibniz rule gives
\cref{eq:rg-thermo-account-jet-product}. Differentiate $ar=1$.
At order $\alpha\ne0$, the term with $\beta=0$ is $a_0r_\alpha$;
all other terms have reciprocal order strictly less than $|\alpha|$.
Solving for $r_\alpha$ gives
\cref{eq:rg-thermo-account-jet-reciprocal}, and induction on total
order proves uniqueness. A local polynomial realization of $a$ stays
positive near its positive base value, so these are the actual derivatives
of its reciprocal. The same observation applies to the finite sum $Z$.
Each active raw mass stays positive in a common neighborhood, and its
normalized mass is $w(x)/Z$. Applying the product rule to that quotient
gives \cref{eq:rg-thermo-account-jet-normalization}. Summing over
$x$ yields the derivatives of $Z/Z=1$, giving exactly the normalization
constraints; zero columns remain identically zero. The argument uses only
scalar reciprocals of positive masses or partitions, never a matrix
inverse. $\square$

\propositionheading{Second-order Fisher jets from third-order masses}
{prop:rg-thermo-account-fisher-jet}
For a probability jet satisfying
\cref{eq:rg-thermo-account-jet-admission}, let $r_{\delta,x}$ be the
reciprocal jet of $\pi(x)$ for each positive atom. For
$\alpha\in\mathcal O_2$ and coordinate unit multi-indices $e_i,e_j$,
\begin{equation}
 \partial^\alpha I_{ij}(\theta_0)
 =\sum_{x:\pi_0(x)>0}
   \sum_{\beta+\gamma+\delta=\alpha}
   \frac{\alpha!}{\beta!\gamma!\delta!}
   \pi_{\beta+e_i}(x)\pi_{\gamma+e_j}(x)r_{\delta,x}.
 \label{eq:rg-thermo-account-fisher-jet}
\end{equation}
Thus third-order mass data determine the full second-order Fisher jet,
including singular zeroth-order tensors.\thermoestablishedstatus


\paragraph{Proof.}
On the local positive support supplied by
\Cref{prop:rg-thermo-account-jet-realization}, the Fisher tensor is
$I_{ij}=\sum_x(\partial_i\pi(x))(\partial_j\pi(x))/\pi(x)$.
Apply the three-factor Leibniz rule to each summand. A derivative of
order at most two of $\partial_i\pi$ needs a mass derivative of
order at most three, giving precisely
\cref{eq:rg-thermo-account-fisher-jet}. All sums are finite and the
reciprocals are smooth on the common positive neighborhood. For any
$v\in\R^d$,
\begin{equation}
 v^\top I(\theta_0)v
 =\sum_{x:\pi_0(x)>0}\pi_0(x)
       \left(\sum_i v_i\frac{\pi_{e_i}(x)}{\pi_0(x)}\right)^2
 \geq0.
 \label{eq:rg-thermo-account-fisher-gram}
\end{equation}
This establishes the zeroth-order sign without requiring rank. The
higher derivative matrices need not be positive semidefinite, and a
singular tensor supplies no Riemannian inverse. $\square$

\propositionheading{Differentiated label, conditional, and channel identities}
{prop:rg-thermo-account-labeled-jets}
Partition $\mathsf X$ by geometry labels. Form
$W_{\gamma,\alpha}=\sum_{x\in\mathsf X_\gamma}\pi_\alpha(x)$.
On each fiber with $W_{\gamma,0}>0$, use scalar normalization to form
$p_\gamma=\pi|_{\mathsf X_\gamma}/W_\gamma$. Then, for
$|\alpha|\leq2$,
\begin{equation}
 \begin{aligned}
 \partial^\alpha I(\pi)
 ={}&\partial^\alpha I(W)\\
 &+\sum_{\gamma:W_{\gamma,0}>0}
   \sum_{\beta\leq\alpha}\binom\alpha\beta
       (\partial^\beta W_\gamma)
       (\partial^{\alpha-\beta}I(p_\gamma)).
 \end{aligned}
 \label{eq:rg-thermo-account-weighted-jet}
\end{equation}
Every derivative is evaluated at $\theta_0$. An unoccupied fiber has
zero weighted contribution and no defined conditional probability or
conditional jet. A fixed finite channel pushes each mass row linearly;
if it preserves labels, the categorical law and all its available
mass and Fisher derivatives are unchanged.\thermoestablishedstatus


\paragraph{Proof.}
The local realization gives positive $W_\gamma$ near $\theta_0$ on
occupied fibers, so its quotient defines the actual conditional family.
The full score is
$\nabla\log W_\gamma+\nabla\log p_\gamma(x)$, and the second
term has conditional mean zero by differentiation of its normalization.
Expanding the score Gram gives
$I(\pi)=I(W)+\sum_\gamma W_\gamma I(p_\gamma)$ as a function
identity. Its ordinary product derivatives give
\cref{eq:rg-thermo-account-weighted-jet}, including derivatives of
$W_\gamma$. If $W_{\gamma,0}=0$, every atom in that fiber has zero
base mass and hence an identically zero realization; dividing by the
fiber mass is neither needed nor defined.

For a fixed channel $C$, finite summation gives
$(\pi C)_\alpha(y)=\sum_x\pi_\alpha(x)C(x,y)$. A target with
zero base mass receives no channel incidence from any positive source
atom, because every term is nonnegative. Source atoms of zero base mass
have zero derivative columns. Thus all target derivative rows at such a
target are zero, and the fixed-support condition persists. If $C$ has
zero entries between different labels, summing its normalized rows over
each target label recovers $W_{\gamma,\alpha}$ for every $\alpha$.
The full Fisher formulas and the fixed-area divisions therefore apply
without dropping any derivative term. $\square$

All these finite algebraic statements hold over real input derivatives.
When every supplied mass, factor derivative, and channel entry is rational,
all displayed products, finite sums, positive scalar reciprocals, and
integer-area divisions remain rational. They then define an exact finite
arithmetic interface in every finite parameter dimension. This does not
assert that arbitrary real smooth generative data have rational values or
rational derivatives. A two-coordinate specialization is an ordinary
restriction of the multi-index formulas, not a restriction of the
mathematical theorem to two parameters.\thermoestablishedstatus


\subsection{Physical native raw-factor and partition jets}
\label{sec:rg-thermo-account-native-jets}

Use the complete native free square of
\Cref{def:rg-pressure-native}. Its edge set $E_L$ contains the two
outgoing signs per cell, and its internal plaquette set
$\mathcal P_L$ has $(L-1)^2$ elements. For
$z\in\{-1,1\}^{E_L}$, put $P_q(z)=\prod_{e\in\partial q}z_e$.
With $|b|<1$, $|t|<1$, and normalized counting reference
$\nu_L(z)=2^{-2L^2}$, define
\begin{equation}
 \begin{aligned}
 R_L(b,t;z)&=\prod_{e\in E_L}(1+bz_e)
                  \prod_{q\in\mathcal P_L}(1+tP_q(z)),\\
 Z_L(b,t)&=2^{-2L^2}\sum_zR_L(b,t;z).
 \end{aligned}
 \label{eq:rg-thermo-account-native-raw}
\end{equation}
The high-temperature setting additionally restricts $b>0$ and
$0<t<1/12$; no volume-limit assertion is used in the finite formulas
below.\thermodeclareddefinition


\propositionheading{Complete native ordinary derivatives}
{prop:rg-thermo-account-native-jet}
For nonnegative integers $r,s$, direct differentiation gives
\begin{align}
 \partial_b^r\partial_t^sR_L
 ={}&r!s!\sum_{\substack{A\subset E_L,\ |A|=r\\
                         D\subset\mathcal P_L,\ |D|=s}}
 \Biggl[\left(\prod_{e\in A}z_e\right)
 \left(\prod_{q\in D}P_q(z)\right)\notag\\
 &\qquad\cdot\prod_{e\in E_L\setminus A}(1+bz_e)
 \prod_{q\in\mathcal P_L\setminus D}(1+tP_q(z))\Biggr],
 \label{eq:rg-thermo-account-native-derivatives}\\
 \partial_b^r\partial_t^s Z_L
 ={}&2^{-2L^2}\sum_z\partial_b^r\partial_t^sR_L.
 \label{eq:rg-thermo-account-native-partition-derivatives}
\end{align}
An empty subset sum is zero and an empty product is one. The normalized
mass jet is the quotient of $\nu_LR_L$ by this same physical $Z_L$.\thermoestablishedstatus


\paragraph{Proof.}
Each edge factor is affine in $b$ and each plaquette factor is affine in
$t$. A nonzero $r$th $b$ derivative must differentiate $r$ distinct edge
factors; every such subset has $r!$ orders of differentiation. The
independent choice of $s$ distinct plaquette factors contributes $s!$.
This gives \cref{eq:rg-thermo-account-native-derivatives}, including
zero when there are too few factors. The finite reference is fixed in
$b,t$, so differentiation commutes with its finite sum and gives
\cref{eq:rg-thermo-account-native-partition-derivatives}. Every factor
is positive on the stated domain, hence $Z_L>0$ and the scalar quotient
recurrence produces the normalized mass derivatives. There are no outer
crossing plaquettes and no fixed exterior in this inventory. $\square$

For $L=1$, there are two independent signs and no plaquette, so the
normalized counting expectation of each edge factor is one and $Z_1=1$.
The $b$ score of a sign is $z/(1+bz)$ and has Fisher information
$1/(1-b^2)$; independence gives $I_{bb}=2/(1-b^2)$, while the
$t$ direction is null. For $L=2$, the edge-factor product with reference
$2^{-8}$ is already a normalized independent-sign law of mean $b$.
Its one internal plaquette has mean $b^4$, so
\begin{equation}
 \begin{aligned}
 Z_2&=1+tb^4,\\
 \partial_b\log p_2&=\partial_b\log R_2-
                  \frac{4tb^3}{1+tb^4},\\
 \partial_t\log p_2&=\partial_t\log R_2-
                  \frac{b^4}{1+tb^4}.
 \end{aligned}
 \label{eq:rg-thermo-account-small-native-partition}
\end{equation}
These examples exhibit the physical normalization terms explicitly. A
finite arithmetic restriction to at most eight independently enumerated
signs permits exactly $L=1,2$ because $2L^2\leq8$; it places no
restriction on the general finite algebraic theorem and supplies no
thermodynamic evidence.

At rational $(b,t)$, the raw and partition derivatives in
\cref{eq:rg-thermo-account-native-derivatives,eq:rg-thermo-account-native-partition-derivatives}
and the resulting normalized Fisher jet are rational. For a different
smooth chart, ordinary composition differentiates every varying source
factor. In particular, the selected all-positive sensor identity of
\Cref{prop:rg-pressure-sensor} multiplies the raw partition by
$(A^{\mathrm{sens}})^{2L^2}$ and evaluates the native factor at its
shifted field $b'$. A hierarchical posterior retains the original
$Z_L(b,t)$ divisor. Derivatives of a varying sensor strength therefore
enter both the scalar and $b'$, as well as the original branch
normalization. For the thermodynamic posterior, the matching interface
is \Cref{prop:rg-thermo-label-representatives}: impose its strictly
positive $C^4$ sensor data and its $C^4$ original and shifted source
maps, both with images in common compact admitted high-temperature
boxes. \Cref{thm:rg-thermo-native-cfour-pressure} then supplies the
actual free-boundary pressure limits at both maps. Their composition,
the $2L^2\log A^{\mathrm{sens}}$ term, and the original-partition
subtraction give the stated $C^3$ effective growth. The selected branch
is the shifted native free law, so
\cref{eq:rg-thermo-account-native-density-inputs} applies to that
shifted map and its exact fixed-word image. The finite source-shift
identity alone supplies none of the source regularity, compact-box,
strict-gap, or center-ratio conditions. An arbitrary positive likelihood
is not covered by this all-edge sensor specialization.\thermoestablishedstatus

\subsection{Prescribed geometry actions and relative label weights}
\label{sec:rg-thermo-account-actions}

An action constant within a branch is distinct from an action common to
all random variables. Fix real geometry-only functions
$\mathscr A_{\gamma,L}(\theta)$, independent of the branch state, and
replace the positive raw geometry weight by
$\widetilde a_{\gamma,L}=a_{\gamma,L}e^{-\mathscr A_{\gamma,L}}$.
For a prescribed normalized hierarchy this means explicitly reweighting
and renormalizing its label masses by the same factor; it is a new
specified ensemble. Conditional branch laws and their physical partitions
remain unchanged. A factor depending on the state would instead change
the branch and needs a separate admission argument.\thermoestablishedstatus


\propositionheading{Action growth and the exact relative multiplier}
{prop:rg-thermo-account-action-growth}
Suppose $\log a_{\gamma,L}$ is $C^3$ on a neighborhood of each
compact $K\subset\Theta$, with
\begin{equation}
 \|\log a_{\gamma,L}\|_{C^3(K)}\leq C_KV,
 \qquad V^{-1}\log a_{\gamma,L}\longrightarrow\alpha_\gamma
                       \quad\text{in }C^3(K).
 \label{eq:rg-thermo-account-raw-growth}
\end{equation}
If the prescribed actions satisfy
\begin{equation}
 \|\mathscr A_{\gamma,L}\|_{C^3(K)}\leq C'_KV,
 \qquad V^{-1}\mathscr A_{\gamma,L}
       \longrightarrow\mathscr A_\gamma^{\mathrm{bulk}}
                       \quad\text{in }C^3(K),
 \label{eq:rg-thermo-account-action-admission}
\end{equation}
then the reweighted raw data have the same growth class, with bulk
log weight $\alpha_\gamma-\mathscr A_\gamma^{\mathrm{bulk}}$.
For either normalization mode, the exact relative label law is
\begin{equation}
 \frac{\widetilde W_{\gamma,L}}{\widetilde W_{\gamma_0,L}}
 =\frac{W_{\gamma,L}}{W_{\gamma_0,L}}
   \exp\bigl[-\mathscr A_{\gamma,L}
                   +\mathscr A_{\gamma_0,L}\bigr].
 \label{eq:rg-thermo-account-relative-action}
\end{equation}\thermoestablishedstatus


\paragraph{Proof.}
The logarithm of the new raw weight is
$\log a_{\gamma,L}-\mathscr A_{\gamma,L}$. Subtraction preserves
the $C^3$ bound of order $V$ and gives the claimed bulk limit.
In the global-factor ensemble the unnormalized label mass is
$a_{\gamma,L}Z_{\gamma,L}$, and in the hierarchy generated from
raw weights it is $a_{\gamma,L}$. Since $Z_{\gamma,L}$ is unchanged,
either log-mass representative is shifted by
$-\mathscr A_{\gamma,L}$. Taking the ratio to $\gamma_0$ cancels
the new common normalization and gives
\cref{eq:rg-thermo-account-relative-action}. The same cancellation
applies when a prescribed hierarchy is explicitly reweighted. All
quantities are finite and positive at fixed volume, so the resulting
finite catalog has a finite positive normalizer. $\square$

An extensive action $V\mathscr A_\gamma(\theta)$ changes the bulk
log-mass density by $-\mathscr A_\gamma$. A fixed smooth order-one action
has zero bulk rate. An action $V^\kappa\mathscr A_\gamma$ with
$0<\kappa<1$ is also subextensive in the full $C^3$ sense, but its
finite value and derivatives remain in
\cref{eq:rg-thermo-account-relative-action}. For two identical fixed
native branches with positive $b,t$ in the high-temperature domain,
initial equal weights and actions
$\mathscr A_{0,L}=0$, $\mathscr A_{1,L}=V^\kappa\theta$ give
$W_{1,L}/W_{0,L}=e^{-V^\kappa\theta}$. At $\theta=0$ the
categorical Fisher information is $V^{2\kappa}/4$, so its extensive
density is $V^{2\kappa-1}/4$. The common native law is independent of
$\theta$ and contributes no score in this example. Thus a vanishing
bulk action does not justify deleting its derivatives: it can vanish
in the extensive density, add a finite correction, or dominate that
density according to whether $\kappa$ is below, equal to, or above
$1/2$. These real geometry actions define positive finite weights even
where their values are negative; nonnegativity is an additional property
of particular energy constructions, not a requirement for this finite
normalization argument.

\propositionheading{Conditional limits and a bounded-action obstruction}
{prop:rg-thermo-account-action-ratios}
Let $\theta_{c,L}$ be declared centers and suppose the original positive
relative weights there converge to $k_\gamma\in(0,\infty)$.
If
\begin{equation}
 \mathscr A_{\gamma,L}(\theta_{c,L})-
 \mathscr A_{\gamma_0,L}(\theta_{c,L})\longrightarrow a_\gamma,
 \label{eq:rg-thermo-account-action-center-input}
\end{equation}
then the new relative weights converge to $k_\gamma e^{-a_\gamma}$.
When these are all labels in the catalog, their normalized weights
converge to the normalization of those relative masses. Boundedness of
the actions, even with the bulk convergence in
\cref{eq:rg-thermo-account-action-admission}, does not imply
\cref{eq:rg-thermo-account-action-center-input} or a limiting ratio.\thermoestablishedstatus


\paragraph{Proof.}
Evaluate \cref{eq:rg-thermo-account-relative-action} at the centers.
Continuity of the exponential and multiplication proves the stated
relative limit. Normalization is continuous on a finite positive vector,
giving the full-catalog assertion. For fixed $C^3$ actions
$\mathscr A_\gamma$ and centers tending to $\theta_0$, continuity
supplies the particular limit
$a_\gamma=\mathscr A_\gamma(\theta_0)-
\mathscr A_{\gamma_0}(\theta_0)$. A common action has $a_\gamma=0$
and cancels exactly at every volume, as the relative formula requires.

For the obstruction, choose two identical fixed interacting native
branches with $0<b<1$ and $0<t<1/12$, and initial unit raw weights.
Their finite physical partitions coincide, so both normalization modes
have equal labels initially. Set
\begin{equation}
 \mathscr A_{0,L}=0,\qquad
 \mathscr A_{1,L}=1+(-1)^L.
 \label{eq:rg-thermo-account-alternating-action}
\end{equation}
These are nonnegative, bounded by two, and independent of $\theta$.
Their parameter derivatives vanish; their $C^3$ norms divided by
$V=L^2$ tend to zero. They satisfy
\cref{eq:rg-thermo-account-action-admission} with zero bulk actions.
Nevertheless, the ratio of label one to label zero is $e^{-2}$ for even
$L$ and one for odd $L$, so the corresponding label-one weights are
$e^{-2}/(1+e^{-2})$ and $1/2$. They do not converge. No phase or
regularity boundary of the native branches has been crossed. The failure
is the missing relative center limit, not a failure of positive finite
normalization. $\square$

\subsection{Which curvature supplies a prescribed action}
\label{sec:rg-thermo-account-curvature-types}

The three curvature objects retain the types fixed in
\Cref{def:rg-geometry-curvature-factor}. Gauge curvature belongs to a
declared principal connection $\omega$; base curvature belongs to a
declared nondegenerate base metric; Fisher curvature belongs to a declared
regular normalized statistical family and its Fisher metric. Naming one
of these objects does not specify the pairing, contraction, integration
measure, coupling, or normalization needed for an action. Once those
data give a finite geometry-only function with
\cref{eq:rg-thermo-account-action-admission}, the preceding propositions
apply to it without identifying the three curvatures.\thermoestablishedstatus


For example, use a unitary representation of dimension $n$, and write
$F_{\mu\nu}$ for its represented, anti-Hermitian gauge curvature.
Let the base carry a prescribed Riemannian metric
$g^{\mathrm{base}}$, raise both base indices with that metric, and use
its volume measure. With $\kappa_{\mathrm g}\geq0$, the sign and
normalization of \cref{eq:rg-geometry-gauge-action} are
\begin{equation}
 \mathscr S_{\mathrm g}
 =-\frac{\kappa_{\mathrm g}}{4n}
   \int_{\mathcal C}
       \operatorname{Re}\Tr(F_{\mu\nu}F^{\mu\nu})
                     d\operatorname{vol}_{g^{\mathrm{base}}}.
 \label{eq:rg-thermo-account-unitary-action}
\end{equation}
The integral must be finite on the prescribed geometry. If this
construction is used for each finite geometry label, its parameter
regularity and extensive growth must still be checked through
\cref{eq:rg-thermo-account-action-admission}; the sign of the
integrand does not establish those bounds.\thermoestablishedstatus

\paragraph{Proof: Pairing and sign.}
For anti-Hermitian represented matrices, the real bilinear pairing
$\langle X,Y\rangle=-\operatorname{Re}\Tr(XY)/n$ satisfies
$\langle X,X\rangle=\Tr(X^*X)/n\geq0$, with equality only for
$X=0$. In an orthonormal Riemannian base frame, the contraction in
\cref{eq:rg-thermo-account-unitary-action} is a sum of these matrix
squares over $\mu,\nu$, proving its nonnegative sign. The corresponding
unitary loop expression also has a direct square form,
\begin{equation}
 \frac{\operatorname{Re}\Tr(2I-U-U^*)}{2n}
 =\frac{\Tr((I-U)^*(I-U))}{2n}\geq0,
 \qquad U^*U=I.
 \label{eq:rg-thermo-account-unitary-loop}
\end{equation}
These signs use both the represented matrix class and the Riemannian
base contraction. They do not assert a positive contraction for an
indefinite base metric. Nor does a nonnegative action alone normalize a
probability measure on all connections; that would require a declared
reference measure and a finite nonzero partition function. $\square$\thermoestablishedstatus


For real $\mathrm{GL}(n)$, $n\geq2$, no positive norm on the entire
matrix Lie algebra is invariant under every conjugation. Indeed,
conjugation of $E_{12}$ by $\operatorname{diag}(s,1,\ldots,1)$,
$s>0$, gives $sE_{12}$. Invariance and norm homogeneity would imply
$\|E_{12}\|=s\|E_{12}\|$ for every $s$, contradicting positivity
on $E_{12}\ne0$. In particular, an unqualified trace-square expression
cannot be interpreted as that positive invariant norm.\thermoestablishedstatus


A separately prescribed SPD fiber metric $H$ supplies a different,
explicitly metric-dependent construction:
\begin{equation}
 \Tr(H^{-1}F^\top HF)
 =\bigl\|H^{1/2}FH^{-1/2}\bigr\|_{\mathrm F}^2\geq0.
 \label{eq:rg-thermo-account-spd-pairing}
\end{equation}
For $G\in\mathrm{GL}(n)$, take the simultaneous changes
$F'=GFG^{-1}$ and $H'=G^{-\top}HG^{-1}$. Substitution into the
left-hand side of \cref{eq:rg-thermo-account-spd-pairing} gives
$\Tr(GH^{-1}F^\top HFG^{-1})$, equal to the original trace by
cyclicity. Positivity therefore comes with the additional covariant
metric data; it is not a norm invariant under arbitrary conjugation of
$F$ while holding all metric data fixed. A base contraction and measure
would still have to be specified before integrating this quantity.\thermoestablishedstatus

A base-curvature functional or a functional of the curvature of a
separately prescribed statistical family can likewise be evaluated as
geometry data, if its metric is regular and the resulting action meets
the finite and $C^3$ growth assumptions above. If instead the action is
required to be the Fisher curvature of the same unknown reweighted law,
the weight equation becomes a new self-consistency problem. The present
normalization argument proves no existence, uniqueness, or regularity
for that problem. It also does not introduce a measure on random base
metrics or connections, an all-depth inverse, a critical or continuum
limit, or an identification with a physical clock or physical law.\thermoestablishedstatus

% Unmatched generic-channel and center assumptions remain explicit.

\endgroup
